% ======================================================================
% Chapter 4: Interacting Fields and Feynman Diagrams
% (book pages 77-130; PDF pages 100-153)
% ======================================================================
\chapter{Interacting Fields and Feynman Diagrams}

\section{Perturbation Theory---Philosophy and Examples}

We have now discussed in some detail the quantization of two free field
theories that give approximate descriptions of many of the particles found in
Nature.  Up to this point, however, free-particle states have been
eigenstates of the Hamiltonian; we have seen no interactions and no
scattering.  In order to obtain a closer description of the real world, we
must include new, nonlinear terms in the Hamiltonian (or Lagrangian) that
will couple different Fourier modes (and the particles that occupy them) to
one another.  To preserve causality, we insist that the new terms may
involve only products of fields at the same spacetime point:
$[\phi(x)]^4$ is fine, but $\phi(x)\phi(y)$ is not allowed.  Thus the terms
describing the interactions will be of the form
\begin{equation}
  H_{\mathrm{int}} = \int d^3x\ \mathcal{H}_{\mathrm{int}}[\phi(x)]
  = -\int d^3x\ \mathcal{L}_{\mathrm{int}}[\phi(x)] .
\end{equation}
For now we restrict ourselves to theories in which
$\mathcal{H}_{\mathrm{int}}$ ($= -\mathcal{L}_{\mathrm{int}}$) is a function
only of the fields, not of their derivatives.

In this chapter we will discuss three important examples of interacting
field theories.  The first is ``phi-fourth'' theory,
\begin{equation}
  \mathcal{L} = \tfrac{1}{2}(\partial_\mu\phi)^2 - \tfrac{1}{2}m^2\phi^2
    - \frac{\lambda}{4!}\phi^4 ,
  \label{eq:4.1}
\end{equation}
where $\lambda$ is a dimensionless coupling constant.  ($\lambda\phi^3$
interaction would be a bit simpler, but then the energy would not be
positive-definite unless we added a higher even power of $\phi$ as well.)
Although we are introducing this theory now for purely pedagogical reasons
(since it is the simplest of all interacting quantum theories), models of the
real world do contain $\phi^4$ interactions; the most important example in
particle physics is the self-interaction of the Higgs field in the standard
electroweak theory.  In Part~II, we will see that $\phi^4$ theory also arises
in statistical mechanics.  The equation of motion for $\phi^4$ theory is
\begin{equation}
  (\partial^2 + m^2)\phi = -\frac{\lambda}{3!}\phi^3,
  \label{eq:4.2}
\end{equation}
which cannot be solved by Fourier analysis as the free Klein--Gordon
equation could.  In the quantum theory we impose the equal-time commutation
relations
\begin{equation}
  [\phi(\mathbf{x}), \pi(\mathbf{y})] = i\delta^{(3)}(\mathbf{x}-\mathbf{y}),
\end{equation}
which are unaffected by $\mathcal{L}_{\mathrm{int}}$.  (Note, however, that
if $\mathcal{L}_{\mathrm{int}}$ contained $\partial_\mu\phi$, the definition
of $\pi(\mathbf{x})$ would change.)  It is an easy exercise to write down the
Hamiltonian of this theory and find the Heisenberg equation of motion for
the operator $\phi(x)$; the result is the same as the classical equation of
motion \eqref{eq:4.2}, just as it was in the free theory.

Our second example of an interacting field theory will be Quantum
Electrodynamics:
\begin{equation}
  \mathcal{L}_{\mathrm{QED}} = \mathcal{L}_{\mathrm{Dirac}}
    + \mathcal{L}_{\mathrm{Maxwell}} + \mathcal{L}_{\mathrm{int}}
  = \bar{\psi}(i\gamma^\mu\partial_\mu - m)\psi
    - \tfrac{1}{4}(F_{\mu\nu})^2 - e\bar{\psi}\gamma^\mu\psi A_\mu ,
  \label{eq:4.3}
\end{equation}
where $A^\mu$ is the electromagnetic vector potential,
$F_{\mu\nu} = \partial_\mu A_\nu - \partial_\nu A_\mu$ is the electromagnetic
field tensor, and $e = -|e|$ is the electron charge.  (To describe a fermion
of charge $Q$, replace $e$ with $Q$.  If we wish to consider several species
of charged particles at once, we simply duplicate $\mathcal{L}_{\mathrm{Dirac}}$
and $\mathcal{L}_{\mathrm{int}}$ for each additional species.)  That such a
simple Lagrangian can account for nearly all observed phenomena from
macroscopic scales down to $10^{-13}\ \mathrm{cm}$ is rather astonishing.
In fact, the QED Lagrangian can be written even more simply:
\begin{equation}
  \mathcal{L}_{\mathrm{QED}} = \bar{\psi}(i\slashed{D} - m)\psi
    - \tfrac{1}{4}(F_{\mu\nu})^2 ,
  \label{eq:4.4}
\end{equation}
where $D_\mu$ is the gauge covariant derivative,
\begin{equation}
  D_\mu = \partial_\mu + ieA_\mu(x) .
  \label{eq:4.5}
\end{equation}
A crucial property of the QED Lagrangian is that it is invariant under the
gauge transformation
\begin{equation}
  \psi(x) \to e^{i\alpha(x)}\psi(x),\qquad
  A_\mu \to A_\mu - \frac{1}{e}\partial_\mu\alpha(x),
  \label{eq:4.6}
\end{equation}
which is realized on the Dirac field as a local phase rotation.  This
invariance under local phase rotations has a fundamental geometrical
significance, which motivates the term covariant derivative.  For our present
purposes, though, it is sufficient just to recognize \eqref{eq:4.6} as a
symmetry of the theory.

The equations of motion follow from \eqref{eq:4.3} by the canonical
procedure.  The Euler--Lagrange equation for $\bar{\psi}$ is
\begin{equation}
  (i\slashed{\partial} - m)\psi(x) = e\slashed{A}\psi(x),
  \label{eq:4.7}
\end{equation}
which is just the Dirac equation coupled to the electromagnetic field by the
minimal coupling prescription, $\partial_\mu \to D_\mu$.  The Euler--Lagrange
equation for $A_\mu$ is
\begin{equation}
  \partial_\mu F^{\mu\nu} = e\bar{\psi}\gamma^\nu\psi = e j^\nu .
  \label{eq:4.8}
\end{equation}
These are the inhomogeneous Maxwell equations, with the current density
$j^\nu = \bar{\psi}\gamma^\nu\psi$ given by the conserved Dirac vector current
\eqref{eq:3.78}.  As with $\phi^4$ theory, the equations of motion can also
be obtained as the Heisenberg equations of motion for the operators $\psi(x)$
and $A^\mu(x)$.  This is easy to verify for $\psi(x)$; we have not yet
discussed the quantization of the electromagnetic field.

In fact, we will not discuss canonical quantization of the electromagnetic
field at all in this book.  It is an awkward subject, essentially because of
gauge invariance.  Note that since $A^0$ does not appear in the Lagrangian
\eqref{eq:4.3}, the momentum conjugate to $A^0$ is identically zero.  This
contradicts the canonical commutation relation $[A^0(\mathbf{x}),
\pi^0(\mathbf{y})] = i\delta(\mathbf{x}-\mathbf{y})$.  One solution is to
quantize in Coulomb gauge, where $\nabla\cdot\mathbf{A} = 0$ and $A^0$ is a
constrained, rather than dynamical, variable; but then manifest Lorentz
invariance is sacrificed.  Alternatively, one can quantize the field in
Lorentz gauge, $\partial_\mu A^\mu = 0$.  It is then possible to modify the
Lagrangian, adding an $A^0$ term.  One obtains the commutation relations
\begin{equation}
  [A^\mu(\mathbf{x}), A^\nu(\mathbf{y})] = -ig^{\mu\nu}\delta(\mathbf{x}-\mathbf{y}),
\end{equation}
essentially the same as four Klein--Gordon fields.  But the extra minus sign
in $[A^0, A^0]$ leads to another (surmountable) difficulty: states created by
$a_{\mathbf{p}}^{0\dagger}$ have negative norm.%
\footnote{Excellent treatments of both quantization procedures are readily
available.  For Coulomb gauge quantization, see Bjorken and Drell (1965),
Chapter~14; for Lorentz gauge quantization, see Mandl and Shaw (1984),
Chapter~5.}

The Feynman rules for calculating scattering amplitudes that involve
photons are derived more easily in the functional integral formulation of
field theory, to be discussed in Chapter~9.  That method has the added
advantage of generalizing readily to the case of non-Abelian gauge fields, as
we will see in Part~III.  In the present chapter we will simply guess the
Feynman rules for photons.  This will actually be quite easy after we derive
the rules for an analogous but simpler theory, Yukawa theory:
\begin{equation}
  \mathcal{L}_{\mathrm{Yukawa}} = \mathcal{L}_{\mathrm{Dirac}}
    + \mathcal{L}_{\mathrm{Klein--Gordon}} - g\bar{\psi}\psi\phi .
  \label{eq:4.9}
\end{equation}
This will be our third example.  It is similar to QED, but with the photon
replaced by a scalar particle $\phi$.  The interaction term contains a
dimensionless coupling constant $g$, analogous to the electron charge $e$.
Yukawa originally invented this theory to describe nucleons ($\psi$) and
pions ($\phi$).  In modern particle theory, the Standard Model contains
Yukawa interaction terms coupling the scalar Higgs field to quarks and
leptons; most of the free parameters in the Standard Model are Yukawa
coupling constants.

Having written down our three paradigm interactions, let us pause a moment
to discuss what other interactions could be found in Nature.  At first it
might seem that the list would be infinite; even for a scalar theory we could
write down interactions of the form $\phi^n$ for any $n$.  But remarkably,
one simple and reasonable axiom eliminates all but a few of the possible
interactions.  That axiom is that the theory be \emph{renormalizable}, and it
arises as follows.  Higher-order terms in perturbation theory, as mentioned
in Chapter~1, will involve integrals over the 4-momenta of intermediate
(``virtual'') particles.  These integrals are often formally divergent, and it
is generally necessary to impose some form of cut-off procedure; the simplest
is just to cut off the integral at some large but finite momentum $\Lambda$.
At the end of the calculation one takes the limit $\Lambda\to\infty$, and
hopes that physical quantities turn out to be independent of $\Lambda$.  If
this is indeed the case, the theory is said to be renormalizable.

Suppose, however, that the theory includes interactions whose coupling
constants have the dimensions of mass to some negative power.  Then to obtain
a dimensionless scattering amplitude, this coupling constant must be
multiplied by some quantity of positive mass dimension, and it turns out that
this quantity is none other than $\Lambda$.  Such a term diverges as
$\Lambda\to\infty$, so the theory is not renormalizable.

We will discuss these matters in detail in Chapter~10.  For now we merely
note that any theory containing a coupling constant with negative mass
dimension is not renormalizable.  A bit of dimensional analysis then allows
us to throw out nearly all candidate interactions.  Since the action
$S = \int\mathcal{L}\,d^4x$ is dimensionless, $\mathcal{L}$ must have
dimension $(\mathrm{mass})^4$ (or simply dimension 4).  From the kinetic terms
of the various free Lagrangians, we note that the scalar and vector fields
$\phi$ and $A^\mu$ have dimension 1, while the spinor field $\psi$ has
dimension $3/2$.  We can now tabulate all of the allowed renormalizable
interactions.

For theories involving only scalars, the allowed interaction terms are
\begin{equation}
  \mu\phi^3 \qquad\text{and}\qquad \lambda\phi^4 .
\end{equation}
The coupling constant $\mu$ has dimension 1, while $\lambda$ is
dimensionless.  Terms of the form $\phi^n$ for $n>4$ are not allowed, since
their coupling constants would have dimension $4-n$.  Of course, more
interesting theories can be obtained by including several scalar fields, real
or complex (see Problem~4.3).

Next we can add spinor fields.  Spinor self-interactions are not allowed,
since $\bar{\psi}\psi^3$ (besides violating Lorentz invariance) already has
dimension $9/2$.  Thus the only allowable new interaction is the Yukawa term,
\begin{equation}
  g\bar{\psi}\psi\phi,
\end{equation}
although similar interactions can also be constructed out of Weyl and
Majorana spinors.

When we add vector fields, many new interactions are possible.  The most
familiar is the vector-spinor interaction of QED,
\begin{equation}
  e\bar{\psi}\gamma^\mu\psi A_\mu .
\end{equation}
Again it is easy to construct similar terms out of Weyl and Majorana
spinors.  Less important is the scalar QED Lagrangian,
\begin{equation}
  \mathcal{L} = |D_\mu\phi|^2 - m^2|\phi|^2,
\end{equation}
which contains
\begin{equation}
  eA^\mu\phi^{*}\partial_\mu\phi,\qquad e^2|\phi|^2A^2 .
\end{equation}
This is our first example of a derivative interaction; quantization of this
theory will be much easier with the functional integral formalism, so we
postpone its discussion until Chapter~9.  Other possible Lorentz-invariant
terms involving vectors are
\begin{equation}
  A^2(\partial_\mu A^\mu) \qquad\text{and}\qquad A^4 .
\end{equation}
Although it is far from obvious, these terms lead to inconsistencies unless
their coupling constants are precisely chosen on the basis of a special type
of symmetry, which must involve several vector fields.  This symmetry
underlies the non-Abelian gauge theories, which will be the main subject of
Part~III.  A mass term $\tfrac{1}{2}m^2A^2$ for vector fields is also
inconsistent, except in the special case where it is added to QED; in any
case, it breaks (Abelian or non-Abelian) gauge invariance.

This exhausts the list of possible Lagrangians involving scalar, spinor, and
vector particles.  It is interesting to note that the currently accepted
models of the strong, weak, and electromagnetic interactions include all of
the types of interactions listed above.  The three paradigm interactions to
be studied in this chapter cover nearly half of the possibilities; we will
study the others in detail later in this book.

The assumption that realistic theories must be renormalizable is certainly
convenient, since a nonrenormalizable theory would have little predictive
power.  However, one might still ask why Nature has been so kind as to use
only renormalizable interactions.  One might have expected that the true
theory of Nature would be a quantum theory of a much more general type.  But
it can be shown that, however complicated a fundamental theory appears at
very high energies, the low-energy approximation to this theory that we see
in experiments should be a renormalizable quantum field theory.  We will
demonstrate this in Section~12.1.

At a more practical level, the preceding analysis highlights a great
difference in methodology between nonrelativistic quantum mechanics and
relativistic quantum field theory.  Since the potential $V(\mathbf{r})$ that
appears in the Schr\"odinger equation is completely arbitrary, nonrelativistic
quantum mechanics puts no limits on what interactions can be found in the
real world.  But we have just seen that quantum field theory imposes very
tight constraints on Nature (or vice versa).  Taken literally, our discussion
implies that the only tasks left for particle physicists are to enumerate the
elementary particles that exist and to measure their masses and coupling
constants.  While this viewpoint is perhaps overly arrogant, the fact that it
is even thinkable is surely a sign that particle physicists are on the right
track toward a fundamental theory.

Given a set of particles and couplings, we must still work out the
experimental consequences.  How do we analyze the quantum mechanics of an
interacting field theory?  It would be nice if we could explicitly solve at
least a few examples (that is, find the exact eigenvalues and eigenvectors as
we did for the free theories) to get a feel for the properties of interacting
theories.  Unfortunately, this is easier said than done.  No exactly solvable
interacting field theories are known in more than two spacetime dimensions,
and even there the solvable models involve special symmetries and
considerable technical complication.%
\footnote{These theories, known as conformal field theories, are the subject
of a large current literature.  See, for example, the reviews by Ginsparg
(1988) and Cardy (1987).}  Studying these theories would be interesting, but
hardly worth the effort at this stage.  Instead we will fall back on a much
simpler and more generally applicable approach: Treat the interaction term
$H_{\mathrm{int}}$ as a perturbation, compute its effects as far in
perturbation theory as is practicable, and hope that the coupling constant is
small enough that this gives a reasonable approximation to the exact answer.
In fact, the perturbation series we obtain will turn out to be very simple in
structure; through the use of Feynman diagrams it will be possible at least
to visualize the effects of interactions to arbitrarily high order.

This simplification of the perturbation series for relativistic field
theories was the great advance of Tomonaga, Schwinger, and Feynman.  To
achieve this simplification, each, independently, found a way to reformulate
quantum mechanics to remove the special role of time, and then applied his
new formalism to electrodynamics.  The interaction picture formalism that we
will use here is that of Tomonaga and Schwinger; Feynman's approach, which
we will take up in Chapter~9, is formulated in terms of the path integral.

\section{Perturbation Expansion of Correlation Functions}

The perturbation-theoretic approach to field theory rests on a few
assumptions and results from quantum mechanics that we will state here in a
way that suits our purposes.  In the quantum mechanics of a particle in a
time-independent potential, one often works with the Schr\"odinger picture,
in which operators are time-independent and the state obeys the Schr\"odinger
equation.  This picture is not well suited to field theory, since the fields
that appear in a Hamiltonian such as \eqref{eq:4.1} have explicit
time-dependent arguments even in the Heisenberg picture.  The
\emph{interaction picture} is a hybrid: one writes
\begin{equation}
  \phi(x) = U^{\dagger}(t,t_0)\ \phi_I(x)\ U(t,t_0),
  \label{eq:4.10}
\end{equation}
where $t$ is the time coordinate of the spacetime point $x$, and
\begin{equation}
  \phi_I(x) = e^{iH_0(t-t_0)}\phi(\mathbf{x},t_0)e^{-iH_0(t-t_0)}
\end{equation}
is the field evolved under the \emph{free} Hamiltonian $H_0$, which is the
interaction Hamiltonian in the Schr\"odinger picture.  The interaction
picture field $\phi_I$ obeys the free-field equations of motion, so it has
the explicit Fourier decomposition we derived in Chapter~2.  The
interaction-picture Hamiltonian is
\begin{equation}
  H_I = e^{iH_0(t-t_0)}H_{\mathrm{int}}e^{-iH_0(t-t_0)} .
  \label{eq:4.11}
\end{equation}

We will also need the following result, which we state without proof: The
interacting-field ground state $|\Omega\rangle$ of the full Hamiltonian $H$
can be expressed in terms of the free vacuum $|0\rangle$ through a
time-ordered exponential.  Concretely, the Heisenberg-picture field $\phi(x)$
is related to the interaction-picture field $\phi_I(x)$ by
\begin{equation}
  \phi(x) = U^{\dagger}(t, t_0)\,\phi_I(x)\,U(t, t_0),
  \label{eq:4.12}
\end{equation}
where $U$ is the time-evolution operator.  If $t_0 = -\infty(1-i\epsilon)$
and $t = +\infty(1-i\epsilon)$, the relation between the two vacuums is
\begin{equation}
  |\Omega\rangle = \lim_{T\to\infty(1-i\epsilon)}
    \bigl( e^{-iE_0T}\langle\Omega|0\rangle \bigr)^{-1}\ e^{-iHT}\ |0\rangle ,
  \label{eq:4.13}
\end{equation}
where $E_0 = \langle\Omega|H|\Omega\rangle$.  The factor $e^{-iE_0T}$ is a
c-number that will cancel against a similar factor in the denominator of our
final expressions.

We now want to compute the two-point correlation function
\begin{equation}
  \langle\Omega|T\phi(x)\phi(y)|\Omega\rangle .
  \label{eq:4.14}
\end{equation}
Using \eqref{eq:4.12}--\eqref{eq:4.13}, this can be rewritten as
\begin{equation}
  \langle\Omega|T\phi(x)\phi(y)|\Omega\rangle
  = \lim_{T\to\infty(1-i\epsilon)}
    \frac{\langle 0|T\bigl\{ \phi_I(x)\phi_I(y)
      \exp\bigl[-i\int_{-T}^{T}dt\ H_I(t)\bigr]\bigr\}|0\rangle}
      {\langle 0|T\bigl\{ \exp\bigl[-i\int_{-T}^{T}dt\ H_I(t)\bigr]\bigr\}|0\rangle} .
  \label{eq:4.15}
\end{equation}
This is the Gell-Mann--Low formula.  Expanding the exponential in powers of
the interaction Hamiltonian gives the perturbation series.

The problem now is to express the full Heisenberg field $\phi$ in terms of
$\phi_I$.  Formally, it is just
\begin{equation}
  \phi(x) = U^{\dagger}(t, t_0)\,\phi_I(x)\,U(t, t_0),
  \label{eq:4.16}
\end{equation}
where we have defined the unitary operator
\begin{equation}
  U(t, t_0) = e^{iH_0(t-t_0)}e^{-iH(t-t_0)},
  \label{eq:4.17}
\end{equation}
known as the interaction picture propagator or time-evolution operator.  We
would like to express $U(t, t_0)$ entirely in terms of $\phi_I$, for which we
have an explicit expression in terms of ladder operators.  To do this, we
note that $U(t, t_0)$ is the unique solution, with initial condition
$U(t_0, t_0) = 1$, of a simple differential equation (the Schr\"odinger
equation):
\begin{equation}
  i\frac{\partial}{\partial t}U(t, t_0)
    = e^{iH_0(t-t_0)}(H-H_0)e^{-iH(t-t_0)}
    = H_I(t)U(t, t_0),
  \label{eq:4.18}
\end{equation}
where
\begin{equation}
  H_I(t) = e^{iH_0(t-t_0)}H_{\mathrm{int}}e^{-iH_0(t-t_0)}
  = \int d^3x\ \mathcal{H}_I(x)
  \label{eq:4.19}
\end{equation}
is the interaction Hamiltonian written in the interaction picture.  The
solution of this differential equation for $U(t, t_0)$ should look something
like $U \sim \exp(-iH_It)$; this would be our desired formula for $U$ in
terms of $\phi_I$.  Doing it more carefully, we will show that the actual
solution is the following power series in $\lambda$:
\begin{equation}
  \begin{aligned}
    U(t, t_0) = 1 &+ (-i)\int_{t_0}^{t}dt_1\ H_I(t_1)
      + (-i)^2\int_{t_0}^{t}dt_1\int_{t_0}^{t_1}dt_2\ H_I(t_1)H_I(t_2)\\
    &+ (-i)^3\int_{t_0}^{t}dt_1\int_{t_0}^{t_1}dt_2\int_{t_0}^{t_2}dt_3\
      H_I(t_1)H_I(t_2)H_I(t_3) + \cdots .
  \end{aligned}
  \label{eq:4.20}
\end{equation}
To verify this, just differentiate: Each term gives the previous one times
$-iH_I(t)$.  The initial condition $U(t_0, t_0) = 1$ for $t=t_0$ is obviously
satisfied.  Note that the various factors of $H_I$ in \eqref{eq:4.20} stand
in time order, later on the left.  This allows us to simplify the expression
considerably, using the time-ordering symbol $T$.  The $H_I^2$ term, for
example, can be written
\begin{equation}
  \int_{t_0}^{t}dt_1\int_{t_0}^{t_1}dt_2\ H_I(t_1)H_I(t_2)
    = \frac{1}{2}\int_{t_0}^{t}dt_1\int_{t_0}^{t}dt_2\
      T\{H_I(t_1)H_I(t_2)\} .
  \label{eq:4.21}
\end{equation}
The double integral on the right-hand side just counts everything twice,
since in the $t_1$--$t_2$ plane, the integrand $T\{H_I(t_1)H_I(t_2)\}$ is
symmetric about the line $t_1 = t_2$ (see Fig.~\ref{fig:4.1}).

\begin{figure}[t]
  \centering
  \includegraphics[width=0.4\textwidth]{fig4.1.png}
  \caption{Geometric interpretation of Eq.~\eqref{eq:4.21}.}
  \label{fig:4.1}
\end{figure}

A similar identity holds for the higher terms:
\begin{equation}
  \int_{t_0}^{t}dt_1\int_{t_0}^{t_1}dt_2\cdots\int_{t_0}^{t_{n-1}}dt_n\
    H_I(t_1)\cdots H_I(t_n)
    = \frac{1}{n!}\int_{t_0}^{t}dt_1\cdots\int_{t_0}^{t}dt_n\
      T\{H_I(t_1)\cdots H_I(t_n)\} .
\end{equation}
This case is a little harder to visualize, but it is not hard to convince
oneself that it is true.  Using this identity, we can now write $U(t, t_0)$
in an extremely compact form:
\begin{equation}
  U(t, t_0) = 1 + (-i)\int_{t_0}^{t}dt_1\ H_I(t_1)
    + \frac{(-i)^2}{2}\int_{t_0}^{t}dt_1\int_{t_0}^{t}dt_2\
      T\{H_I(t_1)H_I(t_2)\} + \cdots
  = T\Bigl\{ \exp\Bigl[ -i\int_{t_0}^{t}dt_1\ H_I(t_1) \Bigr] \Bigr\},
  \label{eq:4.22}
\end{equation}
where the time-ordering of the exponential is just defined as the Taylor
series with each term time-ordered.  When we do real computations we will
keep only the first few terms of the series; the time-ordered exponential is
just a compact way of writing and remembering the correct expression.

We now have control over $\phi(t,\mathbf{x})$; we have written it entirely in
terms of $\phi_I$ as desired.  Before moving on to consider $|\Omega\rangle$,
however, it is convenient to generalize the definition of $U$, allowing its
second argument to take on values other than our ``reference time'' $t_0$.
The correct definition is quite natural:
\begin{equation}
  U(t, t') = T\Bigl\{ \exp\Bigl[ -i\int_{t'}^{t}dt''\ H_I(t'') \Bigr] \Bigr\},
  \qquad (t > t') .
  \label{eq:4.23}
\end{equation}
Several properties follow from this definition, and it is necessary to
verify them.  First, $U(t, t')$ satisfies the same differential equation
\eqref{eq:4.18},
\begin{equation}
  i\frac{\partial}{\partial t}U(t, t') = H_I(t)U(t, t'),
  \label{eq:4.24}
\end{equation}
but now with the initial condition $U = 1$ for $t = t'$.  From this equation
you can show that
\begin{equation}
  U^{\dagger}(t, t') = U^{-1}(t, t'),
  \label{eq:4.25}
\end{equation}
which proves that $U$ is unitary.  Finally, $U(t, t')$ satisfies the
following identities (for $t_1 > t_2 > t_3$):
\begin{equation}
  U(t_1, t_2)U(t_2, t_3) = U(t_1, t_3);
  \qquad U(t_1, t_3)\bigl[U(t_2, t_3)\bigr]^{-1} = U(t_1, t_2) .
  \label{eq:4.26}
\end{equation}

Now we can go on to discuss $|\Omega\rangle$.  Since $|\Omega\rangle$ is the
ground state of $H$, we can isolate it by the following procedure.  Imagine
starting with $|0\rangle$, the ground state of $H_0$, and evolving through
time with $H$:
\begin{equation}
  e^{-iHT}|0\rangle = \sum_n e^{-iE_nT}|n\rangle\langle n|0\rangle,
\end{equation}
where $E_n$ are the eigenvalues of $H$.  We must assume that $|\Omega\rangle$
has some overlap with $|0\rangle$, that is, $\langle\Omega|0\rangle\neq 0$
(if this were not the case, $H_I$ would in no sense be a small perturbation).
Then the above series contains $|\Omega\rangle$, and we can write
\begin{equation}
  e^{-iHT}|0\rangle = e^{-iE_0T}|\Omega\rangle\langle\Omega|0\rangle
    + \sum_{n\neq 0}e^{-iE_nT}|n\rangle\langle n|0\rangle .
\end{equation}
Since $E_n > E_0$ for all $n\neq 0$, we can get rid of all the $n\neq 0$
terms in the series by sending $T$ to $\infty$ in a slightly imaginary
direction: $T\to\infty(1-i\epsilon)$.  Then the exponential factor
$e^{-iE_nT}$ dies slowest for $n=0$, and we have
\begin{equation}
  |\Omega\rangle = \lim_{T\to\infty(1-i\epsilon)}
    \bigl( e^{-iE_0T}\langle\Omega|0\rangle \bigr)^{-1}\ e^{-iHT}\ |0\rangle .
  \label{eq:4.27}
\end{equation}
Since $T$ is now very large, we can shift it by a small constant:
\begin{equation}
  \begin{aligned}
    |\Omega\rangle &= \lim_{T\to\infty(1-i\epsilon)}
      \bigl( e^{-iE_0(t_0-(-T))}\langle\Omega|0\rangle \bigr)^{-1}
      \ e^{-iH(t_0-(-T))}\ |0\rangle\\
    &= \lim_{T\to\infty(1-i\epsilon)}
      \bigl( e^{-iE_0(t_0-(-T))}\langle\Omega|0\rangle \bigr)^{-1}
      \ e^{-iH(t_0-(-T))}e^{-iH_0(-T-t_0)}\ |0\rangle\\
    &= \lim_{T\to\infty(1-i\epsilon)}
      \bigl( e^{-iE_0(t_0-(-T))}\langle\Omega|0\rangle \bigr)^{-1}
      \ U(t_0, -T)\ |0\rangle .
  \end{aligned}
  \label{eq:4.28}
\end{equation}
In the second line we have used $H_0|0\rangle = 0$.  Ignoring the c-number
factor in front, this expression tells us that we can get $|\Omega\rangle$ by
simply evolving $|0\rangle$ from time $-T$ to time $t_0$ with the operator
$U$.  Similarly, we can express $\langle\Omega|$ as
\begin{equation}
  \langle\Omega| = \lim_{T\to\infty(1-i\epsilon)}
    \frac{\langle 0|\,U(T, t_0)}
      {e^{-iE_0(T-t_0)}\langle 0|\Omega\rangle} .
\end{equation}

Now we are ready to evaluate the numerator of \eqref{eq:4.15}.  The most
efficient way to do this is through \emph{Wick's theorem}, which we now
introduce.  Wick's theorem expresses a time-ordered product of fields as a
sum of normal-ordered products with all possible contractions.  We define the
\emph{contraction} of two fields by
\begin{equation}
  \begin{array}{c}
    \overbrace{\phi_I(x)\phi_I(y)}
  \end{array}
  = \begin{cases}
    [\phi_I^+(x), \phi_I^-(y)] & \text{if } x^0 > y^0,\\
    [\phi_I^-(y), \phi_I^+(x)] & \text{if } y^0 > x^0,
  \end{cases}
  = D_F(x-y) ,
  \label{eq:4.29}
\end{equation}
where $\phi_I^+$ and $\phi_I^-$ are the creation and annihilation parts of the
field.  Wick's theorem states that a time-ordered product can be expanded as
\begin{equation}
  T\{\phi(x_1)\cdots\phi(x_m)\} = N\{
    \phi(x_1)\cdots\phi(x_m) + \text{all possible contractions} \},
  \label{eq:4.30}
\end{equation}
where the notation ``all possible contractions'' means that we sum over every
way of contracting pairs of fields (with each contracted pair replaced by its
contraction $D_F$), keeping the fields in their original order.

Let us prove Wick's theorem.  Naturally the proof is by induction on $m$,
the number of fields.  We have already proved the case $m=2$.  So assume the
theorem is true for $m-1$ fields, and let's try to prove it for $m$ fields.
Without loss of generality, we can restrict ourselves to the case
$x^0_1 > x^0_2 > \cdots > x^0_m$; if this is not the case we can just relabel
the points, without affecting either side of \eqref{eq:4.30}.  Then applying
Wick's theorem to $\phi_2\cdots\phi_m$, we have
\begin{equation}
  \begin{aligned}
    T\{\phi_1\cdots\phi_m\} &= \phi_1\cdots\phi_m\\
    &= \phi_1\,N\Bigl\{ \phi_2\cdots\phi_m + \bigl(\text{all contractions}
      \text{ not involving } \phi_1\bigr) \Bigr\}\\
    &= \phi_1^+N\Bigl\{ \cdots \Bigr\}
      + \phi_1^-N\Bigl\{ \cdots \Bigr\} .
  \end{aligned}
  \label{eq:4.31}
\end{equation}
We want to move the $\phi_1^+$ and $\phi_1^-$ inside the $N\{\}$.  For the
$\phi_1^+$ term this is easy: Just move it in, since (being on the left) it is
already in normal order.  The term with $\phi_1^-$ must be put in normal
order by commuting it to the right past all the other $\phi$'s.  Consider,
for example, the term with no contractions:
\begin{equation}
  \begin{aligned}
    \phi_1^-N(\phi_2\cdots\phi_m) &= N(\phi_2\cdots\phi_m)\phi_1^-
      + N\Bigl\{ [\phi_1^-, \phi_2]\phi_3\cdots\phi_m
        + \phi_2[\phi_1^-, \phi_3]\phi_4\cdots\phi_m + \cdots \Bigr\}\\
    &= N\bigl\{ \phi_1^-\phi_2\cdots\phi_m
      + \phi_1\phi_2\phi_3\cdots\phi_m + \cdots \bigr\} .
  \end{aligned}
  \label{eq:4.32}
\end{equation}
The first term in the last line combines with part of the $\phi_1^+$ term
from \eqref{eq:4.31} to give $N\{\phi_1\phi_2\cdots\phi_m\}$, so we now have
the first term on the right-hand side of Wick's theorem, as well as all
possible terms involving a single contraction of $\phi_1$ with another field.
Similarly, a term in \eqref{eq:4.31} involving one contraction will produce
all possible terms involving both that contraction and a contraction of
$\phi_1$ with one of the other fields.  Doing this with all the terms of
\eqref{eq:4.31}, we eventually get all possible contractions of all the
fields, including $\phi_1$.  Thus the induction step is complete, and Wick's
theorem is proved.

\section{Feynman Diagrams}

Wick's theorem allows us to turn any expression of the form
\begin{equation}
  \langle 0|T\{\phi_I(x_1)\phi_I(x_2)\cdots\phi_I(x_n)\}|0\rangle
\end{equation}
into a sum of products of Feynman propagators.  Now we are ready to develop a
diagrammatic interpretation of such expressions.  Consider first the case of
four fields, all at different spacetime points, which we worked out in
Eq.~\eqref{eq:4.30}.  Let us represent each of the points $x_1$ through $x_4$
by a dot, and each factor $D_F(x-y)$ by a line joining $x$ to $y$.  Then
Eq.~\eqref{eq:4.30} can be represented as the sum of three diagrams (called
Feynman diagrams):
\begin{equation}
  \langle 0|T\phi_1\phi_2\phi_3\phi_4|0\rangle = D_F(x_1-x_2)D_F(x_3-x_4)
    + D_F(x_1-x_3)D_F(x_2-x_4)
    + D_F(x_1-x_4)D_F(x_2-x_3) .
  \label{eq:4.42}
\end{equation}
Although this isn't exactly a measurable quantity, the diagrams do suggest an
interpretation: Particles are created at two spacetime points, each propagates
to one of the other points, and then they are annihilated.  This can happen in
three ways, corresponding to the three ways to connect the points in pairs, as
shown in the three diagrams.  The total amplitude for the process is the sum
of the three diagrams.

Things get more interesting when the expression contains more than one field
at the same spacetime point.  So let us now return to the evaluation of the
two-point function $\langle\Omega|T\{\phi(x)\phi(y)\}|\Omega\rangle$, and put
formula \eqref{eq:4.15} to use.  We will ignore the denominator until the
very end of this section.  The numerator, with the exponential expanded as a
power series, is
\begin{equation}
  \langle 0|T\Bigl\{ \phi(x)\phi(y)
    \Bigl[ 1 + (-i)\int dt\ H_I(t) + \cdots \Bigr] \Bigr\}|0\rangle .
  \label{eq:4.43}
\end{equation}
The first term gives the free-field result,
$\langle 0|T\{\phi(x)\phi(y)\}|0\rangle = D_F(x-y)$.  The second term, in
$\phi^4$ theory, is
\begin{equation}
  \langle 0|T\Bigl\{ \phi(x)\phi(y)(-i)\int dt\int d^3z\
    \frac{\lambda}{4!}\phi^4 \Bigr\}|0\rangle
  = \langle 0|T\Bigl\{ \phi(x)\phi(y)
    \Bigl( -\frac{i\lambda}{4!} \Bigr)\int d^4z\
      \phi(z)\phi(z)\phi(z)\phi(z) \Bigr\}|0\rangle .
\end{equation}
Now apply Wick's theorem.  We get one term for every way of contracting the
six $\phi$ operators with each other in pairs.  There are 15 ways to do this,
but (fortunately) only two of them are really different.  If we contract
$\phi(x)$ with $\phi(y)$, then there are three ways to contract the four
$\phi(z)$'s with each other, and all three give identical expressions.  The
other possibility is to contract $\phi(x)$ with one of the $\phi(z)$'s (four
choices), $\phi(y)$ with one of the others (three choices), and the remaining
two $\phi(z)$'s with each other (one choice).  There are twelve ways to do
this, and all give identical expressions.  Thus we have
\begin{equation}
  \begin{aligned}
    \langle 0|T\Bigl\{ \phi(x)\phi(y)(-i)\int dt\int d^3z\
      \frac{\lambda}{4!}\phi^4 \Bigr\}|0\rangle
    &= 3\cdot\Bigl( -\frac{i\lambda}{4!} \Bigr)\ D_F(x-y)\
      \int d^4z\ D_F(z-z)D_F(z-z)\\
    &\quad + 12\cdot\Bigl( -\frac{i\lambda}{4!} \Bigr)\
      \int d^4z\ D_F(x-z)D_F(y-z)D_F(z-z) .
  \end{aligned}
  \label{eq:4.44}
\end{equation}
We can understand this expression better if we represent each term as a
Feynman diagram.  Again we draw each contraction $D_F$ as a line, and each
point as a dot.  But this time we must distinguish between the ``external''
points, $x$ and $y$, and the ``internal'' point $z$; each internal point is
associated with a factor of $(-i\lambda)\int d^4z$.  We will worry about the
constant factors in a moment.  Using these rules, we see that the above
expression \eqref{eq:4.44} is equal to the sum of two diagrams:
\begin{equation}
  \begin{aligned}
    &\underbrace{---\!(x)---\!\bigcirc\!---(y)---}_{\text{disconnected}}
      \qquad + \qquad
    \underbrace{(x)\!-\!-\!\cdot\!-\!-\!-\!\bigcirc\!-\!-\!-\!\cdot\!-\!-\!(y)}
      {\text{connected}} .
  \end{aligned}
\end{equation}
We refer to the lines in these diagrams as \emph{propagators}, since they
represent the propagation amplitude $D_F$.  Internal points, where four lines
meet, are called \emph{vertices}.  Since $D_F(x-y)$ is the amplitude for a
free Klein--Gordon particle to propagate between $x$ and $y$, the diagrams
actually interpret the analytic formula as a process of particle creation,
propagation, and annihilation which takes place in spacetime.

Now let's try a more complicated contraction, from the $\lambda^3$ term in
the expansion of the correlation function:
\begin{equation}
  \begin{aligned}
    &\langle 0|\ \phi(x)\phi(y)\ \frac{1}{3!}\Bigl( -\frac{i\lambda}{4!}
      \Bigr)^3 \int d^4z\,d^4w\,d^4u\ \phi^4(z)\phi^4(w)\phi^4(u)\ |0\rangle\\
    &= \int d^4z\,d^4w\,d^4u\ D_F(x-z)D_F(z-z)D_F(z-w)
      \times D_F(w-y)D_F^2(w-u)D_F(u-u) .
  \end{aligned}
  \label{eq:4.45}
\end{equation}
The number of ``different'' contractions that give this same expression is
large:
\begin{equation}
  \underbrace{3!}_{\substack{\text{interchange}\\ \text{of vertices}}}
  \times \underbrace{\frac{4\cdot 3}{2}}_{\substack{\text{placement of}\\
    \text{contractions into }z}}
  \times \underbrace{\frac{4\cdot 3\cdot 2}{2!}}_{\substack{\text{placement}
    \text{ of contractions into }w}}
  \times \underbrace{\frac{4\cdot 3\cdot 2}{2!}}_{\substack{\text{placement}
    \text{ of contractions into }u}}
  \times \underbrace{\frac{1}{2}}_{\substack{\text{interchange of}\\
    w-u \text{ contractions}}} .
\end{equation}
The product of these combinatoric factors is 10,368, roughly $1/13$ of the
total of 135,135 possible full contractions of the 14 operators.  The
structure of this particular contraction can be represented by the following
``cactus'' diagram:
\begin{equation}
  (x)\!-\!-\!\bigcirc_{z}\!-\!-\!\bigcirc_{w}\!-\!-\!(y)
  \qquad\text{with self-contractions on } z \text{ and } u .
\end{equation}
It is conventional, for obvious reasons, to let this one diagram represent
the sum of all 10,368 identical terms.

In practice one always draws the diagram first, using it as a mnemonic device
for writing down the analytic expression.  But then the question arises, What
is the overall constant?  We could, of course, work it out as above: We could
associate a factor $\int d^4z\,(-i\lambda/4!)$ with each vertex, put in the
$1/n!$ from the Taylor series, and then do the combinatorics by writing out
the product of fields as in \eqref{eq:4.45} and counting.  But the $1/n!$ from
the Taylor series will almost always cancel the $n!$ from interchanging the
vertices, so we can just forget about both of these factors.  Furthermore,
the generic vertex has four lines coming in from four different places, so the
various placements of these contractions into $\phi\phi\phi\phi$ generates a
factor of $4!$ (as in the $w$ vertex above), which cancels the denominator in
$(-i\lambda/4!)$.  It is therefore conventional to associate the expression
$\int d^4z\,(-i\lambda)$ with each vertex.  (This was the reason for the
factor of $4!$ in the $\phi^4$ coupling.)

In the above diagram, this scheme gives a constant that is too large by a
factor of $8 = 2\cdot2\cdot2$, the \emph{symmetry factor} of the diagram.
Two factors of 2 come from lines that start and end on the same vertex: The
diagram is symmetric under the interchange of the two ends of such a line.
The other factor of 2 comes from the two propagators connecting $w$ to $u$:
The diagram is symmetric under the interchange of these two lines with each
other.  A third possible type of symmetry is the equivalence of two vertices.
To get the correct overall constant for a diagram, we divide by its symmetry
factor, which is in general the number of ways of interchanging components
without changing the diagram.

Most people never need to evaluate a diagram with a symmetry factor greater
than 2, so there's no need to worry too much about these technicalities.
But here are a few examples, to make some sense out of the above rules:
\begin{equation}
  \begin{array}{ccc}
    \text{diagram } (x)\!-\!-\!(y) & \text{self-loop at } x \text{ and } y &
      \text{two propagators between } x \text{ and } y\\
    S = 2 & S = 2\cdot2\cdot2 = 8 & S = 2
  \end{array}
\end{equation}
When in doubt, you can always determine the symmetry factor by counting
equivalent contractions, as we did above.

We are now ready to summarize our rules for calculating the numerator
\begin{equation}
  \lim_{T\to\infty(1-i\epsilon)}
    \langle 0|T\Bigl\{ \phi_I(x)\phi_I(y)\exp\Bigl[ -i\int_{-T}^{T}dt\ H_I(t)
      \Bigr] \Bigr\}|0\rangle .
\end{equation}
Draw all possible diagrams with two external points, following these rules:
\begin{enumerate}
  \item For each propagator,
    \begin{equation}
      \begin{array}{c}
        \overbrace{\phi_I(x)\phi_I(y)}
      \end{array} = D_F(x-y),
    \end{equation}
    draw a line connecting $x$ to $y$.
  \item For each vertex,
    \begin{equation}
      (-i\lambda)\int d^4z,
    \end{equation}
    draw a point $z$ with four lines emerging from it.
  \item For each external point, draw a dot.
  \item Divide by the symmetry factor of the diagram.
\end{enumerate}

The explicit dependence on $T$ seems to disappear when we take the limit
$T\to\infty$ in the previous equation.  But consider the diagram
\begin{equation}
  (x)\!-\!-\!\bigcirc
\end{equation}
The delta function for the left-hand vertex is $(2\pi)^4\delta^{(4)}(p_1+p_2)$,
so momentum conservation at the right-hand vertex is automatically satisfied,
and we get $(2\pi)^4\delta^{(4)}(0)$ there.  This awkward factor is easy to
understand by going back to position space.  It is simply the integral of a
constant over $d^4w$:
\begin{equation}
  \int d^4w\ (\text{const}) \propto (2T)\cdot(\text{volume of space})
  \label{eq:4.49}
\end{equation}
This just tells us that the spacetime process Eq.~(4.48) can happen at
any place in space, and at any time between $-T$ and $T$.  Every disconnected
piece of a diagram, that is, every piece that is not connected to an external
point, will have one such $(2\pi)^4\delta^{(4)}(0) = 2TV$ factor.

The contributions to the correlation function coming from such diagrams can
be better understood with the help of a very pretty identity, the
\emph{exponentiation of the disconnected diagrams}.  It works as follows.  A
typical diagram has the form
\begin{equation}
  (x)\!-\!\text{(connected)}\!-\!(y) \quad\text{plus several disconnected
    pieces} .
  \label{eq:4.50}
\end{equation}
with a piece connected to $x$ and $y$, and several disconnected pieces.
(Since each vertex has an even number of lines coming into it, $x$ and $y$
must be connected to each other.)  Label the various possible disconnected
pieces by $V_i$:
\begin{equation}
  V_i \equiv \text{value of disconnected piece } \mathcal{V}_i .
  \label{eq:4.51}
\end{equation}
The elements $\mathcal{V}_i$ are connected internally, but disconnected from
external points.  Suppose that a diagram (such as \eqref{eq:4.50}) has $n_i$
pieces of the form $\mathcal{V}_i$, for each $i$, in addition to its one
piece that is connected to $x$ and $y$.  (In any given diagram, only finitely
many of the $n_i$ will be nonzero.)  If we also let $V_i$ denote the value of
the piece $\mathcal{V}_i$, then the value of such a diagram is
\begin{equation}
  (\text{value of connected piece}) \cdot \prod_i \frac{V_i^{n_i}}{n_i!} .
\end{equation}
The $1/n_i!$ is the symmetry factor coming from interchanging the $n_i$ copies
of $\mathcal{V}_i$.  The sum of all diagrams, representing the numerator of
our formula for the two-point correlation function, is then
\begin{equation}
  \begin{aligned}
    \Bigl( \sum_{\text{connected}} \Bigr)
      &\times \sum_{\text{all } \{n_i\}} \prod_i \frac{V_i^{n_i}}{n_i!}\\
    &= \Bigl( \sum_{\text{connected}} \Bigr)
      \times \prod_i \exp(V_i)
    = \Bigl( \sum_{\text{connected}} \Bigr)
      \times \exp\Bigl( \sum_i V_i \Bigr) .
  \end{aligned}
  \label{eq:4.52}
\end{equation}
where ``all $\{n_i\}$'' means ``all ordered sets $\{n_1, n_2, n_3, \ldots\}$
of nonnegative integers.''  The sum of the connected pieces factors out of
this expression.  We have just shown that the sum of all diagrams is equal to
the sum of all connected diagrams, times the exponential of the sum of all
disconnected diagrams.  (We should really say ``pieces'' rather than
``diagrams'' on the right-hand side of the equality, but from now on we will
often just call a single piece a ``diagram.'')  Pictorially, the identity is
\begin{equation}
  \lim_{T\to\infty(1-i\epsilon)}
    \langle 0|T\Bigl\{ \phi_I(x)\phi_I(y)\exp\Bigl[ -i\int_{-T}^{T}dt\ H_I(t)
      \Bigr]\Bigr\}|0\rangle
  = \Bigl( \sum_{\text{connected}} \Bigr)\times
    \exp\Bigl( \sum_i V_i \Bigr) .
  \label{eq:4.53}
\end{equation}

Now consider the denominator of our formula \eqref{eq:4.15} for the two-point
function.  By an argument identical to the above, it is just
\begin{equation}
  \langle 0|T\Bigl\{ \exp\Bigl[ -i\int_{-T}^{T}dt\ H_I(t) \Bigr]
    \Bigr\}|0\rangle = \exp\Bigl( \sum_i V_i \Bigr) ,
\end{equation}
which cancels the exponential of the disconnected diagrams in the numerator.
This is the final simplification of the formula, which now reads
\begin{equation}
  \langle\Omega|T\{\phi(x)\phi(y)\}|\Omega\rangle
  = \text{sum of all connected diagrams with two external points} .
  \label{eq:4.54}
\end{equation}
We have come a long way from our original formula, Eq.~\eqref{eq:4.15}.

Having gotten rid of the disconnected diagrams in our formula for the
correlation function, we might pause a moment to go back and interpret them
physically.  The place to look is Eq.~\eqref{eq:4.30}, which can be written
\begin{equation}
  \lim_{T\to\infty(1-i\epsilon)}
    \frac{\langle 0|T\bigl\{ \phi_I(x)\phi_I(y)
      \exp[-i\int_{-T}^{T}dt\,H_I(t)]\bigr\}|0\rangle}
      {\langle 0|\Omega\rangle e^{-iE_0(2T)}\langle\Omega|0\rangle}
  = \langle\Omega|T\phi(x)\phi(y)|\Omega\rangle .
\end{equation}
Looking only at the $T$-dependent parts of both sides, this implies
\begin{equation}
  \exp\Bigl[ \sum_i V_i \Bigr] \propto \exp\bigl\{ -iE_0(2T) \bigr\} .
  \label{eq:4.55}
\end{equation}
Since each disconnected diagram $V_i$ contains a factor of
$(2\pi)^4\delta^{(4)}(0) = 2TV$, this gives us a formula for the energy
density of the vacuum (relative to the zero of energy set by
$H_0|0\rangle = 0$):
\begin{equation}
  \frac{E_0}{V} = -\frac{i}{2T}\sum_i V_i
    \Big|_{\text{factor out } (2\pi)^4\delta^{(4)}(0)}
  \label{eq:4.56}
\end{equation}
We should emphasize that the right-hand side is independent of $T$ and
(volume); in particular it is reassuring to see that $E_0$ is proportional to
the volume of space.  In Chapter~11 we will find that this formula is
actually useful.

This completes our present analysis of the two-point correlation function.
The generalization to higher correlation functions is easy:
\begin{equation}
  \langle\Omega|T\phi(x_1)\cdots\phi(x_n)|\Omega\rangle
  = \text{sum of all connected diagrams with } n \text{ external points} .
  \label{eq:4.57}
\end{equation}
The disconnected diagrams exponentiate, factor, and cancel as before, by the
same argument.  There is now a potential confusion in terminology, however.
By ``disconnected'' we mean ``disconnected from all external points''---exactly
the same diagrams as in \eqref{eq:4.51}.  (They are sometimes called
``vacuum bubbles'' or ``vacuum to vacuum transitions''.)  In higher
correlation functions, diagrams can also be disconnected in another sense.
Consider, for example, the four-point function:
\begin{equation}
  \langle\Omega|T\phi_1\phi_2\phi_3\phi_4|\Omega\rangle
  = \text{sum of all diagrams with 4 external points},
  \label{eq:4.58}
\end{equation}
including diagrams where external points are disconnected from each other.
Such diagrams do not exponentiate or factor; they contribute to the amplitude
just as do the fully connected diagrams (in which any point can be reached
from any other by traveling along the lines).

Note that in $\phi^4$ theory, all correlation functions of an odd number of
fields vanish, since it is impossible to draw an allowed diagram with an odd
number of external points.  We can also see this by going back to Wick's
theorem: The interaction Hamiltonian $H_I$ contains an even number of fields,
so all terms in the perturbation expansion of an odd correlation function
will contain an odd number of fields.  But it is impossible to fully contract
an odd number of fields in pairs, and only fully contracted terms have
nonvanishing vacuum expectation value.

\section{Cross Sections and the $S$-Matrix}

We now have an extremely beautiful formula, Eq.~\eqref{eq:4.57}, for
computing an extremely abstract quantity: the $n$-point correlation function.
Our next task is to find equally beautiful ways of computing quantities that
can actually be measured: cross sections and decay rates.  In this section,
after briefly reviewing the definitions of these objects, we will relate them
(via a rather technical but fairly careful derivation) to a more primitive
quantity, the $S$-matrix.  In the next section we will learn how to compute
the matrix elements of the $S$-matrix using Feynman diagrams.

\subsection{The Cross Section}

The experiments that probe the behavior of elementary particles, especially
at high energies, can be divided into two categories: scattering experiments
and decay experiments.  In a scattering experiment, we collide particles with
well-defined momenta and measure the momenta and angular distribution of the
final-state particles.  The central observable is the \emph{cross section}
$\sigma$, which is essentially the probability of a particular scattering
process, normalized to unit incident flux and unit target density.  For a
two-body initial state $A+B$, the cross section is defined by
\begin{equation}
  \sigma = \frac{1}{N_A N_B}\ \times\ \text{(number of scattering events per
    unit time per unit volume)} ,
  \label{eq:4.59}
\end{equation}
where $N_A$ and $N_B$ are the number densities of particles $A$ and $B$.

To compute $\sigma$ in field theory, we need to relate it to the transition
amplitudes of the theory.  The physical picture is as follows.  We imagine
that at time $t\to-\infty$ the system is in a state of free particles, well
separated from each other, which we call the ``in'' state.  As time
increases, the particles interact, and at $t\to+\infty$ the system evolves
into a superposition of free-particle ``out'' states.  The transition
amplitudes between in and out states define the $S$-matrix.  To compute them
in perturbation theory, we use the interaction picture: the in and out states
are related by the $S$-matrix,
\begin{equation}
  \langle \text{out}|S|\text{in}\rangle .
  \label{eq:4.60}
\end{equation}
The $S$-matrix is defined by
\begin{equation}
  S = \lim_{T\to\infty} U(T, -T)
    = T\Bigl\{ \exp\Bigl[ -i\int_{-\infty}^{\infty}dt\ H_I(t) \Bigr] \Bigr\} .
  \label{eq:4.61}
\end{equation}
In the interaction picture, the initial state is a superposition of
free-particle states of definite momentum.  To compute the scattering
amplitude, we consider a wavepacket initial state built from the states
$|\mathbf{k}_A\rangle$ and $|\mathbf{k}_B\rangle$:
\begin{equation}
  |\phi_A\phi_B\rangle_{\mathrm{in}} = \int
    \frac{d^3k_A}{(2\pi)^3}\ \frac{\phi_A(\mathbf{k}_A)e^{-ib\cdot\mathbf{k}_A}}
      {\sqrt{2E_A}}
    \int \frac{d^3k_B}{(2\pi)^3}\ \frac{\phi_B(\mathbf{k}_B)}
      {\sqrt{2E_B}}\ |\mathbf{k}_A\mathbf{k}_B\rangle_{\mathrm{in}} .
  \label{eq:4.68}
\end{equation}
where $b$ is the impact parameter.  We could expand
$\langle\phi_1\phi_2\cdots|$ in terms of similarly defined out states of
definite momentum formed in the asymptotic future.%
\footnote{Here and below, the product symbol applies (symbolically) to the
integral as well as the other factors in parentheses; the integrals apply to
what is outside the parentheses as well.}  It is much easier, however, to use
the out states of definite momentum as the final states in the probability
amplitude, and to multiply by the various normalization factors after
squaring the amplitude.  This is physically reasonable as long as the
detectors of final-state particles mainly measure momentum---that is, they do
not resolve positions at the level of de Broglie wavelengths.

We can now relate the probability of scattering in a real experiment to an
idealized set of transition amplitudes between the asymptotically defined in
and out states of definite momentum,
\begin{equation}
  {}_{\mathrm{out}}\langle \mathbf{p}_1\mathbf{p}_2\cdots
    |\mathbf{k}_A\mathbf{k}_B\rangle_{\mathrm{in}} .
  \label{eq:4.69}
\end{equation}
To compute the overlap of in states with out states, we note that the
conventions for defining the two sets of states are related by time
translation:
\begin{equation}
  \begin{aligned}
    {}_{\mathrm{out}}\langle \mathbf{p}_1\mathbf{p}_2\cdots
      |\mathbf{k}_A\mathbf{k}_B\rangle_{\mathrm{in}}
    &= \lim_{T\to\infty}\ \langle \mathbf{p}_1\mathbf{p}_2\cdots|
      \ e^{-iH(2T)}\ |\mathbf{k}_A\mathbf{k}_B\rangle
    = \lim_{T\to\infty}\ \langle \mathbf{p}_1\mathbf{p}_2\cdots|
      \ U(-T,T)\ |\mathbf{k}_A\mathbf{k}_B\rangle .
  \end{aligned}
  \label{eq:4.70}
\end{equation}
In the last line, the states are defined at any common reference time.  Thus,
the in and out states are related by the limit of a sequence of unitary
operators.  This limiting unitary operator is called the $S$-matrix:
\begin{equation}
  {}_{\mathrm{out}}\langle \mathbf{p}_1\mathbf{p}_2\cdots
    |\mathbf{k}_A\mathbf{k}_B\rangle_{\mathrm{in}}
  \equiv \langle \mathbf{p}_1\mathbf{p}_2\cdots|S|\mathbf{k}_A\mathbf{k}_B\rangle .
  \label{eq:4.71}
\end{equation}
The $S$-matrix has the following structure: If the particles in question do
not interact at all, $S$ is simply the identity operator.  Even if the theory
contains interactions, the particles have some probability of simply missing
one another.  To isolate the interesting part of the $S$-matrix---that is,
the part due to interactions---we define the $T$-matrix by
\begin{equation}
  S = 1 + iT .
  \label{eq:4.72}
\end{equation}
Next we note that the matrix elements of $S$ should reflect 4-momentum
conservation.  Thus $S$ or $T$ should always contain a factor
$(2\pi)^4\delta^{(4)}(\sum p)$.  Extracting this factor, we define the
invariant matrix element $\mathcal{M}$, by
\begin{equation}
  \langle \mathbf{p}_1\mathbf{p}_2\cdots|iT|\mathbf{k}_A\mathbf{k}_B\rangle
    = (2\pi)^4\delta^{(4)}(k_A+k_B-\textstyle\sum p_f)
      \cdot i\mathcal{M}(k_A, k_B \to \{p_f\}) .
  \label{eq:4.73}
\end{equation}
We have written this expression in terms of 4-momenta $p$ and $k$, but of
course all 4-momenta are on mass-shell: $p^0 = E_{\mathbf{p}}$,
$k^0 = E_{\mathbf{k}}$.  (Note that our entire treatment is specific to the
case where the initial state contains only two particles.  For $3\to$many or
many$\to$many interactions, one can invent analogous constructions, but we
will not consider such complicated experiments in this book.)

The matrix element $\mathcal{M}$ is analogous to the scattering amplitude $f$
of one-particle quantum mechanics.  It is useful because it allows us to
separate all the physics that depends on the details of the interaction
Hamiltonian (``dynamics'') from all the physics that doesn't
(``kinematics'').  In the next section we will discuss how to compute
$\mathcal{M}$ using Feynman diagrams.  But first, we must figure out how to
reconstruct the cross section $\sigma$ from $\mathcal{M}$.

To do this, let us calculate, in terms of $\mathcal{M}$, the probability for
the initial state $|\phi_A\phi_B\rangle$ to scatter and become a final state
of $n$ particles whose momenta lie in a small region
$d^3p_1\cdots d^3p_n$.  In our normalization, this probability is
\begin{equation}
  \mathcal{P}(AB \to 12\cdots n)
    = \Bigl( \prod_f \frac{d^3p_f}{(2\pi)^3}\frac{1}{2E_f} \Bigr)
      \Bigl| {}_{\mathrm{out}}\langle \mathbf{p}_1\cdots\mathbf{p}_n
        |\phi_A\phi_B\rangle_{\mathrm{in}} \Bigr|^2 .
  \label{eq:4.74}
\end{equation}
For a single target ($A$) particle and many incident ($B$) particles with
different impact parameters $b$, the number of scattering events is
\begin{equation}
  N = \sum_{\text{all incident particles } i}
    \mathcal{P}_i = \int d^2b\ n_B\,\mathcal{P}(b),
\end{equation}
where $n_B$ is the number density (particles per unit area) of $B$ particles.
Since we are assuming that this number density is constant over the range of
the interaction, $n_B$ can be taken outside the integral.  The cross section
is then
\begin{equation}
  \sigma = \frac{N}{n_BN_A} = \frac{n_B}{n_B}\int d^2b\ \mathcal{P}(b)
    = \int d^2b\ \mathcal{P}(b) .
  \label{eq:4.75}
\end{equation}
Deriving a simple expression for $\sigma$ in terms of $\mathcal{M}$ is now a
fairly straightforward calculation.  Combining \eqref{eq:4.75},
\eqref{eq:4.74}, and \eqref{eq:4.68}, we have (writing $d\sigma$ rather than
$\sigma$ since this is an infinitesimal quantity)
\begin{equation}
  \begin{aligned}
    d\sigma &= \int d^2b\ \prod_f \frac{d^3p_f}{(2\pi)^3}\frac{1}{2E_f}\
      \prod_{i=A,B}\frac{d^3k_i}{(2\pi)^3}\ \frac{\phi_i(\mathbf{k}_i)
        e^{-ib\cdot\mathbf{k}_i}}{\sqrt{2E_i}}\\
    &\quad \times e^{ib\cdot(\mathbf{k}_A'-\mathbf{k}_B')}\
      \bigl( {}_{\mathrm{out}}\langle\{p_f\}|\{k_i\}\rangle_{\mathrm{in}} \bigr)
      \bigl( {}_{\mathrm{out}}\langle\{p_f\}|\{k_i'\}\rangle_{\mathrm{in}}
        \bigr)^{*} ,
  \end{aligned}
  \label{eq:4.76}
\end{equation}
where we have used $k_A'$ and $k_B'$ as dummy integration variables in the
second half of the squared amplitude.  The $d^2b$ integral can be performed
to give a factor of $(2\pi)^2\delta^{(2)}(\mathbf{k}_B-\mathbf{k}_B')$.  We
get more delta functions by writing the final two factors of \eqref{eq:4.76}
in terms of $\mathcal{M}$.  Assuming that we are not interested in the trivial
case of forward scattering where no interaction takes place, we can drop the
1 in Eq.~\eqref{eq:4.72} and write these factors as
\begin{align}
  {}_{\mathrm{out}}\langle\{p_f\}|\{k_i\}\rangle_{\mathrm{in}}
    &= i\mathcal{M}(k_A, k_B \to \{p_f\})\,
      (2\pi)^4\delta^{(4)}(\textstyle\sum k - \sum p_f);\\
  \bigl( {}_{\mathrm{out}}\langle\{p_f\}|\{k_i'\}\rangle_{\mathrm{in}} \bigr)^{*}
    &= -i\mathcal{M}^{*}(k_A', k_B' \to \{p_f\})\,
      (2\pi)^4\delta^{(4)}(\textstyle\sum k' - \sum p_f) .
\end{align}
We can use the second of these delta functions, together with the
$(2\pi)^2\delta^{(2)}(\mathbf{k}_B-\mathbf{k}_B')$, to perform all six of the
$k$ integrals in \eqref{eq:4.76}.  Of the six integrals, only those over
$k_A^z$ and $k_B^z$ require some work:
\begin{equation}
  \int dk_A^z\,dk_B^z\ \delta(k_A^z+k_B^z-\textstyle\sum p_f^z)\,
    \delta(E_A+E_B-\textstyle\sum E_f)
  = \int dk_A^z\ \delta\Bigl(
    \sqrt{k_A^2+m_A^2} + \sqrt{k_B^2+m_B^2} - \textstyle\sum E_f \Bigr)
    \frac{E_B}{k_B^z}\ \frac{E_A}{k_A^z}
    \frac{1}{|v_A - v_B|} .
  \label{eq:4.77}
\end{equation}
In the last line and in the rest of Eq.~\eqref{eq:4.76} it is understood that
the constraints $k_A^z+k_B^z = \sum p_f^z$ and $E_A+E_B = \sum E_f$ now apply
(in addition to the constraints $k_A = k_A'$ and $k_B = k_B'$ coming from the
other four integrals).  The difference $|v_A - v_B|$ is the relative velocity
of the beams as viewed from the laboratory frame.

Now recall that the initial wavepackets are localized in momentum space,
centered on $\mathbf{p}_A$ and $\mathbf{p}_B$.  This means that we can
evaluate all factors that are smooth functions of $k_A$ and $k_B$ at
$\mathbf{p}_A$ and $\mathbf{p}_B$, pulling them outside the integrals.  These
factors include $E_AE_B$, $|v_A-v_B|$, and $\mathcal{M}$---everything except
the remaining delta function.  After doing this, we arrive at the expression
\begin{equation}
  \begin{aligned}
    d\sigma &= \frac{1}{2E_A\,2E_B\,|v_A-v_B|}\
      \prod_f \frac{d^3p_f}{(2\pi)^3}\frac{1}{2E_f}\
      |\mathcal{M}(p_A,p_B\to\{p_f\})|^2\\
    &\quad \times \int \frac{d^3k_A}{(2\pi)^3}\ \frac{d^3k_B}{(2\pi)^3}\
      |\phi_A(\mathbf{k}_A)|^2|\phi_B(\mathbf{k}_B)|^2\
      (2\pi)^4\delta^{(4)}(k_A+k_B-\textstyle\sum p_f) .
  \end{aligned}
  \label{eq:4.78}
\end{equation}
To simplify this formula further, we should think a bit more about the
properties of real particle detectors.  We have already noted that real
detectors project mainly onto eigenstates of momentum.  But real detectors
have finite resolution; that is, they sum incoherently over momentum bites of
finite size.  Normally, the measurement of the final-state momentum is not of
such high quality that it can resolve the small variation of this momentum
that results from the momentum spread of the initial wavepackets
$\phi_A$, $\phi_B$.  In that case, we may treat even the momentum vector
$k_A+k_B$ inside the delta function as being well approximated by its central
value $p_A+p_B$.  With this further approximation, we can perform the
integrals over $k_A$ and $k_B$ using the normalization condition
Eq.~(4.66).  This produces the final form of the relation between
$S$-matrix elements and cross sections,
\begin{equation}
  d\sigma = \frac{1}{2E_A\,2E_B\,|v_A-v_B|}\
    \Bigl( \prod_f \frac{d^3p_f}{(2\pi)^3}\frac{1}{2E_f} \Bigr)\
    |\mathcal{M}(p_A,p_B\to\{p_f\})|^2\
    (2\pi)^4\delta^{(4)}\bigl( p_A+p_B-\textstyle\sum p_f \bigr) .
  \label{eq:4.79}
\end{equation}
All dependence on the shapes of the wavepackets has disappeared.

The integral over final-state momenta in \eqref{eq:4.79} has the structure
\begin{equation}
  \int d\Pi_n = \Bigl( \prod_f \int \frac{d^3p_f}{(2\pi)^3}\frac{1}{2E_f}
    \Bigr)\ (2\pi)^4\delta^{(4)}\bigl( p_A+p_B-\textstyle\sum p_f \bigr),
  \label{eq:4.80}
\end{equation}
with $P$ the total initial 4-momentum.  This integral is manifestly Lorentz
invariant, since it is built up from invariant 3-momentum integrals
constrained by a 4-momentum delta function.  This integral is known as
relativistically invariant $n$-body phase space.  Of the other ingredients in
\eqref{eq:4.79}, the matrix element $\mathcal{M}$ is also Lorentz invariant.
The Lorentz transformation property of \eqref{eq:4.79} therefore comes
entirely from the prefactor
\begin{equation}
  \frac{1}{E_AE_B|v_A-v_B|} .
  \label{eq:4.81}
\end{equation}
This factor is \emph{not} Lorentz invariant.  But the cross section itself
should be invariant under boosts along the beam direction.  The resolution of
this apparent paradox is that the cross section is not a Lorentz scalar
quantity in general; what is invariant is the quantity $E_AE_B|v_A-v_B|\,
d\sigma$.  In practice, one works in a frame where the cross section is
simply related to the invariant matrix element, most often the
center-of-mass frame.  In the CM frame, with both particles massless or with
equal masses, one has $E_A = E_B = E_{\mathrm{cm}}/2$ and
$|v_A-v_B| = 2$, and one obtains the familiar formula
\begin{equation}
  \Bigl( \frac{d\sigma}{d\Omega} \Bigr)_{\mathrm{CM}}
    = \frac{1}{64\pi^2E_{\mathrm{cm}}^2}\ |\mathcal{M}|^2 ,
  \label{eq:4.82}
\end{equation}
which reduces to Eq.~\eqref{eq:1.1} in the massless case.

\section{Computing $S$-Matrix Elements from Feynman Diagrams}

We now wish to evaluate the matrix elements
\begin{equation}
  \langle \mathbf{p}_1\cdots\mathbf{p}_n|S|\mathbf{k}_A\mathbf{k}_B\rangle
  \label{eq:4.83}
\end{equation}
of the $S$-matrix in terms of Feynman diagrams.  The most straightforward
way to do this is to go back to the definition \eqref{eq:4.61} of $S$ and
use Wick's theorem to evaluate the time-ordered products of interaction
Hamiltonians, with the additional step of contracting the fields with the
creation and annihilation operators of the external states.

It is conventional to work with the $T$-matrix element
$\langle\{p_f\}|iT|\{k_i\}\rangle$ and the invariant amplitude
$\mathcal{M}$ defined in \eqref{eq:4.73}.  To compute $\mathcal{M}$, we
evaluate the correlation function
\begin{equation}
  \langle \Omega|T\phi(x_1)\cdots\phi(x_n)|\Omega\rangle ,
\end{equation}
extract the external propagators that carry the external momenta, and amputate
them; equivalently, we use directly the momentum-space Feynman rules derived
below.

For example, consider $\phi^4$ theory.  The lowest-order contribution to the
scattering of two particles $A, B \to 1, 2$ is obtained from the term
\begin{equation}
  \langle \mathbf{p}_1\mathbf{p}_2|S|\mathbf{p}_A\mathbf{p}_B\rangle
  \supset \langle \mathbf{p}_1\mathbf{p}_2|T\Bigl\{ \frac{(-i\lambda)}{4!}
    \int d^4x\ \phi_I^4(x) \Bigr\}|\mathbf{p}_A\mathbf{p}_B\rangle .
  \label{eq:4.91}
\end{equation}
We use Wick's theorem to reduce the product of fields to normal-ordered
form, and then contract the remaining fields with the external states.  The
next term in $\langle \mathbf{p}_1\mathbf{p}_2|S|\mathbf{p}_A\mathbf{p}_B\rangle$
is
\begin{equation}
  \begin{aligned}
    \langle \mathbf{p}_1\mathbf{p}_2|&T\Bigl\{ \frac{(-i\lambda)}{4!}
      \int d^4x\ \phi_I(x) \Bigr\}|\mathbf{p}_A\mathbf{p}_B\rangle_0\\
    &= \langle \mathbf{p}_1\mathbf{p}_2|\frac{(-i\lambda)}{4!}
      \int d^4x\ \bigl( \phi_I(x) + \text{contractions} \bigr)
      |\mathbf{p}_A\mathbf{p}_B\rangle_0 ,
  \end{aligned}
  \label{eq:4.92}
\end{equation}
using Wick's theorem.  Since the external states are not $|0\rangle$, terms
that are not fully contracted do not necessarily vanish; we can use an
annihilation operator from $\phi_I(x)$ to annihilate an initial-state
particle, or a creation operator from $\phi_I(x)$ to produce a final-state
particle.  For example,
\begin{equation}
  \phi_I(x)|\mathbf{p}\rangle = e^{-ip\cdot x}|0\rangle .
  \label{eq:4.93}
\end{equation}
An uncontracted $\phi_I$ operator inside the $T$-product of \eqref{eq:4.92}
has two terms: $\phi_I^+$ on the far right and $\phi_I^-$ on the far left.
We get one contribution to the $S$-matrix element for each way of commuting
the $a$ of $\phi_I^-$ past an initial-state $a^\dagger$, and one contribution
for each way of commuting the $a^\dagger$ of $\phi_I^+$ past a final-state
$a$.  It is natural, then, to define the contractions of field operators with
external states as follows:
\begin{equation}
  \begin{aligned}
    \phi_I(x)|\mathbf{p}\rangle &= e^{-ip\cdot x}|0\rangle;\qquad
    \langle\mathbf{p}|\phi_I(x) = \langle 0|e^{+ip\cdot x} .
  \end{aligned}
  \label{eq:4.94}
\end{equation}
To evaluate an $S$-matrix element such as \eqref{eq:4.92}, we simply write
down all possible full contractions of the $\phi_I$ operators and the
external-state momenta.

To see that this prescription is correct, let us evaluate \eqref{eq:4.92} in
detail.  The $T$-product contains terms of the form
\begin{equation}
  \phi\phi\phi\phi;\qquad \phi\phi\phi\phi \text{ with one contraction};
  \qquad \phi\phi\phi\phi \text{ with two contractions} .
  \label{eq:4.95}
\end{equation}
The last term, in which the $\phi$ operators are fully contracted with each
other, is equal to a vacuum bubble diagram times the value of \eqref{eq:4.91}
calculated above:
\begin{equation}
  \frac{(-i\lambda)}{4!}\int d^4x\ \langle \mathbf{p}_1\mathbf{p}_2|
    \phi^4|\mathbf{p}_A\mathbf{p}_B\rangle_0\ \Bigl|_{\text{fully contracted}}
  = (\text{vacuum bubble})\times \langle \mathbf{p}_1\mathbf{p}_2|
    S|\mathbf{p}_A\mathbf{p}_B\rangle_0 .
  \label{eq:4.96}
\end{equation}
This is just another contribution to the trivial part of the $S$-matrix, so
we ignore it.

Next consider the second term of \eqref{eq:4.95}, in which two of the four
$\phi$ operators are contracted.  The normal-ordered product of the remaining
two fields looks like $(a^\dagger a^\dagger + 2a^\dagger a + aa)$.  As we
commute these operators past the $a$'s and $a^\dagger$'s of the initial and
final states, we find that only a term with an equal number of $a$'s and
$a^\dagger$'s can survive.  In the language of contractions, this says that
one of the $\phi$'s must be contracted with an initial-state $|\mathbf{p}\rangle$,
the other with a final-state $\langle\mathbf{p}|$.  The uncontracted
$|\mathbf{p}\rangle$ and $\langle\mathbf{p}|$ give a delta function as in
\eqref{eq:4.91}.  To represent these quantities diagrammatically, we introduce
external lines to our Feynman rules:
\begin{equation}
  \begin{aligned}
    \phi_I(x)|\mathbf{p}\rangle &=
      \underbrace{-\!-\!-\!-\!-\!\bullet}_{p}\; ;
    \qquad
    \langle\mathbf{p}|\phi_I(x) =
      \underbrace{\bullet\!-\!-\!-\!-\!-\!}_{p}\; .
  \end{aligned}
  \label{eq:4.97}
\end{equation}
Feynman diagrams for $S$-matrix elements will always contain external lines,
rather than the external points of diagrams for correlation functions.  The
second term of \eqref{eq:4.95} thus yields four diagrams, corresponding to
the four ways of assigning the two uncontracted operators to the two external
lines.  The integration $\int d^4x$ produces a momentum-conserving delta
function at each vertex (including the external momenta), so these diagrams
again describe trivial processes in which the initial and final states are
identical.  This illustrates a general principle: Only fully connected
diagrams, in which all external lines are connected to each other, contribute
to the $T$-matrix.

Finally, consider the term of \eqref{eq:4.95} in which none of the operators
are contracted with each other.  Our prescription tells us to contract two of
the $\phi$'s with $|\mathbf{p}_A\mathbf{p}_B\rangle$ and the other two with
$\langle \mathbf{p}_1\mathbf{p}_2|$.  There are $4!$ ways to do this.  Thus we
obtain the diagram
\begin{equation}
  \mathcal{M} \supset (4!)\cdot \frac{(-i\lambda)}{4!}\
    \int d^4x\ e^{-i(p_A+p_B-p_1-p_2)\cdot x}
  = -i\lambda\,(2\pi)^4\delta^{(4)}(p_A+p_B-p_1-p_2) .
  \label{eq:4.98}
\end{equation}
This is exactly of the form $i\mathcal{M}(2\pi)^4\delta^{(4)}(p_A+p_B-p_1-p_2)$,
with $\mathcal{M} = -\lambda$.

Before continuing our discussion of Feynman diagrams for $S$-matrix elements,
we should certainly pause to turn this result into a cross section.  For
scattering in the center-of-mass frame, we can simply plug
$|\mathcal{M}|^2 = \lambda^2$ into Eq.~\eqref{eq:4.82} to obtain
\begin{equation}
  \Bigl( \frac{d\sigma}{d\Omega} \Bigr)_{\mathrm{CM}}
    = \frac{\lambda^2}{64\pi^2E_{\mathrm{cm}}^2} .
  \label{eq:4.99}
\end{equation}
We have just computed our first quantum field theory cross section.  It is a
rather dull result, having no angular dependence at all.  (This situation
will be remedied when we consider fermions in the next section.)  Integrating
over $d\Omega$, and dividing by 2 since there are two identical particles in
the final state, we find the total cross section,
\begin{equation}
  \sigma_{\mathrm{total}} = \frac{\lambda^2}{32\pi E_{\mathrm{cm}}^2} .
  \label{eq:4.100}
\end{equation}
In practice, one would probably use this result to measure the value of
$\lambda$.

Returning to our general discussion, let us consider some higher-order
contributions to the $T$-matrix for the process $A, B \to 1, 2$.  If we
ignore, for the moment, the ``connected and amputated'' prescription, we have
the formula
\begin{equation}
  i\mathcal{M} \sim \text{(lowest order)} + \text{one-loop diagrams}
    + \text{two-loop diagrams} + \cdots ,
  \label{eq:4.101}
\end{equation}
plus diagrams in which the four external lines are not all connected to each
other.  We have already seen that this last class of diagrams gives no
contribution to the $T$-matrix.  The first diagram shown in \eqref{eq:4.101}
gives the lowest-order contribution to $T$, which we calculated above.  The
next three diagrams give expected corrections to this amplitude, involving
creation and annihilation of additional ``virtual'' particles.

The diagrams in the second line of \eqref{eq:4.101} contain disconnected
``vacuum bubbles''.  By the same argument as at the end of Section~4.4, the
disconnected pieces exponentiate to an overall phase factor giving the shift
of the energy of the interacting vacuum state upon which the scattering takes
place.  Thus they are irrelevant to $S$.  We have now seen that only fully
connected diagrams give sensible contributions to $S$-matrix elements.

The last diagram is more problematical; let us evaluate it.  After
integrating over the two vertex positions, we obtain
\begin{equation}
  \begin{aligned}
    \mathcal{M} &\supset \frac{(-i\lambda)^2}{2}\
      \int \frac{d^4p'}{(2\pi)^4}\ \frac{i}{p'^2-m^2}\
      \int \frac{d^4k}{(2\pi)^4}\ \frac{i}{k^2-m^2}\\
    &\quad \times (-i\lambda)(2\pi)^4\delta^{(4)}(p_A+p'-p_1-p_2)\
      (-i\lambda)(2\pi)^4\delta^{(4)}(p_B-p') .
  \end{aligned}
  \label{eq:4.102}
\end{equation}
We can integrate over $p'$ using the second delta function.  It tells us to
evaluate
\begin{equation}
  \frac{i}{p'^2-m^2}\Big|_{p'=p_B} = \frac{i}{p_B^2-m^2} .
\end{equation}
We get infinity, since $p_B$, being the momentum of an external particle, is
on-shell: $p_B^2 = m^2$.  This is a disaster.  Clearly, our formula for $S$
makes sense only if we exclude diagrams of this form, that is, diagrams with
loops connected to only one external leg.  Fortunately, this is physically
reasonable: In the same way that the vacuum bubble diagrams represent the
evolution of $|0\rangle$ into $|\Omega\rangle$, these external leg
corrections,
\begin{equation}
  \underbrace{-\!-\!\bigcirc\!-\!-}_{p},
\end{equation}
represent the evolution of $|\mathbf{p}\rangle_0$ into $|\mathbf{p}\rangle$,
the single-particle state of the interacting theory.  Since these corrections
have nothing to do with the scattering process, we should exclude them from
the computation of $S$.

For a general diagram with external legs, we define \emph{amputation} in the
following way.  Starting from the tip of each external leg, find the last
point at which the diagram can be cut by removing a single propagator, such
that this operation separates the leg from the rest of the diagram.  Cut
there.  For example:
\begin{equation}
  \text{cut the external legs at their last ``attachment point'' to the
    rest of the diagram} .
\end{equation}

Let us summarize our prescription for calculating scattering amplitudes.
Our formula for $S$-matrix elements, Eq.~(4.90), can be rewritten
\begin{equation}
  i\mathcal{M}\cdot(2\pi)^4\delta^{(4)}(p_A+p_B-\textstyle\sum p_f)
    = \sum \text{all connected, amputated diagrams} .
  \label{eq:4.103}
\end{equation}
By ``connected'', we now mean fully connected, that is, with no vacuum
bubbles, and all external legs connected to each other.  The Feynman rules
for scattering amplitudes in $\phi^4$ theory are, in position space,
\begin{enumerate}
  \item For each propagator,
    \begin{equation}
      x\!-\!-\!y = D_F(x-y);
    \end{equation}
  \item For each vertex, a factor $-i\lambda\int d^4z$;
  \item For each external line, the wavefunction factor $e^{-ip\cdot x}$ (or
    $e^{ip\cdot x}$);
  \item Divide by the symmetry factor.
\end{enumerate}
Notice that the factor for an ingoing line is just the amplitude for that
particle to be found at the vertex it connects to, i.e., the particle's
wavefunction.  Similarly, the factor for an outgoing line is the amplitude
for a particle produced at the vertex to have the desired final momentum.

Just as with the Feynman rules for correlation functions, it is usually
simpler to introduce the momentum-space representation of the propagators,
carry out the vertex integrals to obtain momentum-conserving delta functions,
and use these delta functions to evaluate as many momentum integrals as
possible.  In a scattering amplitude, however, there will always be an
overall delta function, which can be used to cancel the one on the left-hand
side of Eq.~\eqref{eq:4.103}.  We are then left with
\begin{equation}
  i\mathcal{M} = \text{sum of all connected, amputated diagrams},
  \label{eq:4.104}
\end{equation}
where the diagrams are evaluated according to the following rules:
\begin{enumerate}
  \item For each propagator,
    \begin{equation}
      \frac{i}{p^2-m^2+i\epsilon} ;
    \end{equation}
  \item For each vertex, $-i\lambda$;
  \item For each external line, a factor of 1;
  \item Impose momentum conservation at each vertex;
  \item Integrate over each undetermined loop momentum $\int d^4k/(2\pi)^4$;
  \item Divide by the symmetry factor.
\end{enumerate}

\section{Feynman Rules for Fermions}

The Feynman rules for fermions are more complicated than those for scalar
fields, because the fermion fields anticommute.  We will develop these rules
in the context of Yukawa theory, which couples a Dirac fermion to a real
scalar field:
\begin{equation}
  \mathcal{L} = \bar{\psi}(i\slashed{\partial} - m)\psi
    + \tfrac{1}{2}(\partial_\mu\phi)^2 - \tfrac{1}{2}m^2\phi^2
    - g\bar{\psi}\psi\phi .
  \label{eq:4.112}
\end{equation}
This is a simplified model of Quantum Electrodynamics.  In this section we
will carefully work out the rules of calculation for Yukawa theory, so that
in the next section we can guess the rules for QED without too much
difficulty.

To be even more specific, consider the two-particle scattering reaction
\begin{equation}
  \text{fermion}(p) + \text{fermion}(k) \to \text{fermion}(p')
    + \text{fermion}(k') .
  \label{eq:4.113}
\end{equation}
The leading contribution comes from the $H_I^2$ term of the $S$-matrix:
\begin{equation}
  \langle \mathbf{p}',\mathbf{k}'|T\Bigl\{ (-ig)\int d^4x\
    \bar{\psi}_I\psi_I\phi_I\ (-ig)\int d^4y\
    \bar{\psi}_I\psi_I\phi_I \Bigr\}|\mathbf{p},\mathbf{k}\rangle_0 .
  \label{eq:4.114}
\end{equation}
To evaluate this expression, use Wick's theorem to reduce the $T$-product to
an $N$-product of contractions, then act the uncontracted fields on the
initial- and final-state particles.  Represent this latter process as the
contraction
\begin{equation}
  \begin{aligned}
    \psi_I(x)|\mathbf{p}, s\rangle &=
      \int \frac{d^3p'}{(2\pi)^3}\ \frac{1}{\sqrt{2E_{\mathbf{p}'}}}\
        \sum_{s'} a_{\mathbf{p}'}^{s'}\,u^{s'}(p')e^{-ip'\cdot x}\
      \sqrt{2E_{\mathbf{p}}}\,a_{\mathbf{p}}^{s\dagger}|0\rangle
    = e^{-ip\cdot x}u^s(p)|0\rangle .
  \end{aligned}
  \label{eq:4.115}
\end{equation}
Similar expressions hold for the contraction of $\psi_I$ with a final-state
fermion, and for contractions of $\psi_I$ and $\bar{\psi}_I$ with
antifermion states.  Note that $\psi_I$ can be contracted with a fermion on
the right or an antifermion on the left; the opposite is true for
$\bar{\psi}_I$.

We can write a typical contribution to the matrix element \eqref{eq:4.114} as
the contraction
\begin{equation}
  \langle \mathbf{p}',\mathbf{k}'|\ \frac{(-ig)^2}{2}\
    \int d^4x\,d^4y\ \bar{\psi}\psi\phi\ \bar{\psi}\psi\phi\
    |\mathbf{p},\mathbf{k}\rangle .
  \label{eq:4.115b}
\end{equation}
Up to a possible minus sign, the value of this quantity is
\begin{equation}
  (-ig)^2\ \int \frac{d^4q}{(2\pi)^4}\ \frac{i}{q^2-m^2}\
    (2\pi)^4\delta^{(4)}(p'-p-q)\ (2\pi)^4\delta^{(4)}(k'-k+q)
    \ \bar{u}(p')u(p)\,\bar{u}(k')u(k) .
  \label{eq:4.115c}
\end{equation}
(We have dropped the factor $1/2!$ because there is a second, identical term
that comes from interchanging $x$ and $y$ in \eqref{eq:4.115b}.)  Using
either delta function to perform the integral, we find that this expression
takes the form $i\mathcal{M}(2\pi)^4\delta^{(4)}(\sum p)$, with
\begin{equation}
  i\mathcal{M} = \frac{-ig^2}{q^2-m^2}\,
    \bar{u}(p')u(p)\,\bar{u}(k')u(k) .
  \label{eq:4.116}
\end{equation}
When writing it in this way, we must remember to impose the constraints
$p - p' = q = k' - k$.

Instead of working from \eqref{eq:4.115b}, we could draw a Feynman diagram:
We denote scalar particles by dashed lines, and fermions by solid lines.  The
$S$-matrix element could then be obtained directly from the following
momentum-space Feynman rules.
\begin{enumerate}
  \item Propagators:
    \begin{align}
      \phi(x)\phi(y) &= \frac{i}{q^2-m^2+i\epsilon};\label{eq:4.117a}\\
      \psi(x)\bar{\psi}(y) &= \frac{i(\slashed{p}+m)}{p^2-m^2+i\epsilon}
        .\label{eq:4.117b}
    \end{align}
  \item Vertices: $-ig$.
  \item External leg contractions:
    \begin{align}
      \psi(x)|p, s\rangle &= u^s(p);\qquad
      \langle p, s|\bar{\psi}(x) = \bar{u}^s(p);\nonumber\\
      \bar{\psi}(x)|k, s\rangle &= \bar{v}^s(k);\qquad
      \langle k, s|\psi(x) = v^s(k) .
    \end{align}
  \item Impose momentum conservation at each vertex.
  \item Integrate over each undetermined loop momentum.
  \item Figure out the overall sign of the diagram.
\end{enumerate}
Several comments are in order regarding these rules.

First, note that the $1/n!$ from the Taylor series of the time-ordered
exponential is always canceled by the $n!$ ways of interchanging vertices to
obtain the same contraction.  The diagrams of Yukawa theory never have
symmetry factors, since the three fields $(\psi\bar{\psi}\phi)$ in $H_I$
cannot substitute for one another in contractions.

Second, the direction of the momentum on a fermion line is always
significant.  On external lines, as for bosons, the direction of the momentum
is always ingoing for initial-state particles and outgoing for final-state
particles.  This follows immediately from the expansions of $\psi$ and
$\bar{\psi}$, where the annihilation operators $a_{\mathbf{p}}$ and
$b_{\mathbf{p}}$ both multiply $e^{-ip\cdot x}$ and the creation operators
$a_{\mathbf{p}}^\dagger$ and $b_{\mathbf{p}}^\dagger$ both multiply
$e^{+ip\cdot x}$.  On internal fermion lines (propagators), the momentum must
be assigned in the direction of particle-number flow (for electrons, this is
the direction of negative charge flow).  This requirement is most easily seen
by working out an example from first principles.  Consider the annihilation
of a fermion and an antifermion into two bosons:
\begin{equation}
  \langle \mathbf{k}\mathbf{k}'|S|\mathbf{p}\mathbf{p}'\rangle
    = \int d^4x\int d^4y\ e^{ik'\cdot x}\ \bar{v}(p')\,e^{-ip'\cdot x}\
      \Bigl[ \frac{i}{q^2-m^2} \Bigr]\ u(p)e^{-ip\cdot y}\ e^{ik\cdot y} .
  \label{eq:4.118}
\end{equation}
The integrals over $x$ and $y$ give delta functions that force $q$ to flow
from $y$ to $x$, as shown.  On internal boson lines the direction of the
momentum is irrelevant and may be chosen for convenience, since
$D_F(x-y) = D_F(y-x)$.

It is conventional to draw arrows on fermion lines, as shown, to represent the
direction of particle-number flow.  The momentum assigned to a fermion
propagator then flows in the direction of this arrow.  For external
antiparticles, however, the momentum flows opposite to the arrow; it helps to
show this explicitly by drawing a second arrow next to the line.

Third, note that in our examples the Dirac indices contract together along
the fermion lines.  This will also happen in more complicated diagrams:
\begin{equation}
  \bar{u}(p_3)\ \frac{i(\slashed{p}_2+m)}{p_2^2-m^2}\
    \frac{i(\slashed{p}_1+m)}{p_1^2-m^2}\ u(p_0) ,
  \label{eq:4.119}
\end{equation}
with the Dirac indices contracted in the order the fields appear along the
line.

Finally, let's take a moment to worry about fermion minus signs.  Return to
the example of the fermion-fermion scattering process.  We adopt a sign
convention for the initial and final states:
\begin{equation}
  |\mathbf{p}\mathbf{k}\rangle \equiv a_{\mathbf{p}}^\dagger
    a_{\mathbf{k}}^\dagger|0\rangle;\qquad
  \langle \mathbf{p}',\mathbf{k}'| \equiv \langle 0|a_{\mathbf{k}'}
    a_{\mathbf{p}'} ,
  \label{eq:4.120}
\end{equation}
so that $(\langle\mathbf{p},\mathbf{k}|)^\dagger = |\mathbf{p},\mathbf{k}\rangle$.
Then the contraction
\begin{equation}
  \langle \mathbf{p}',\mathbf{k}'|(\bar{\psi}\psi)_x(\bar{\psi}\psi)_y
    |\mathbf{p},\mathbf{k}\rangle
  = \langle 0|a_{\mathbf{k}'}a_{\mathbf{p}'}\ a^\dagger a\ a^\dagger a\
    a_{\mathbf{p}}^\dagger a_{\mathbf{k}}^\dagger|0\rangle
\end{equation}
can be untangled by moving $a_{\mathbf{p}'}$ two spaces to the left, and so
picks up a factor of $(-1)^2 = +1$.  But note that in the contraction
\begin{equation}
  \langle 0|a_{\mathbf{k}'}a_{\mathbf{p}'}\ a_{\mathbf{p}}^\dagger a\ a^\dagger
    a_{\mathbf{k}}^\dagger|0\rangle ,
\end{equation}
it is sufficient to move the $\bar{\psi}_y$ one space to the left, giving a
factor of $-1$.  This contraction corresponds to the diagram in which the
two fermion lines are crossed.

The full result, to lowest order, for the $S$-matrix element for this process
is therefore
\begin{equation}
  i\mathcal{M} = \frac{-ig^2}{(p-p')^2-m^2}\,
    \bar{u}(p')u(p)\,\bar{u}(k')u(k)
    - \frac{-ig^2}{(p-k')^2-m^2}\,
    \bar{u}(p')u(k)\,\bar{u}(k')u(p) .
  \label{eq:4.121}
\end{equation}
The minus sign difference between these diagrams is a reflection of Fermi
statistics.  Turning this expression into an explicit cross section would
require some additional work; we postpone such calculations until Chapter~5,
when we can work with QED instead of the less interesting Yukawa theory.

In complicated diagrams, one can often simplify the determination of the
minus signs by noting that the product $(\bar{\psi}\psi)$, or any other pair
of fermions, commutes with any operator.  Thus, the anticommutation relations
are only needed to determine the relative signs of terms in which the order
of these fermion pairs is interchanged.

\section{Feynman Rules for Quantum Electrodynamics}

We are now ready to write down the Feynman rules for QED.  The interaction
Hamiltonian is
\begin{equation}
  H_{\mathrm{int}} = \int d^3x\ e\bar{\psi}\gamma^\mu\psi A_\mu ,
  \label{eq:4.129}
\end{equation}
so the vertex factor is $-ie\gamma^\mu$.  The propagators are as follows.  For
the fermions, we have already derived the Feynman propagator
\begin{equation}
  S_F(p) = \frac{i(\slashed{p}+m)}{p^2-m^2+i\epsilon}
    = \frac{i}{\slashed{p}-m+i\epsilon} .
  \label{eq:4.130}
\end{equation}
For the photon, the external line factors are obtained from the expansion of
the free electromagnetic field,
\begin{equation}
  A^\mu(x) = \int \frac{d^3p}{(2\pi)^3}\ \frac{1}{\sqrt{2E_{\mathbf{p}}}}
    \sum_{r=0}^{3} \Bigl( \epsilon^{r\mu}(p)\,a_{\mathbf{p}}^r e^{-ip\cdot x}
      + \epsilon^{r\mu*}(p)\,a_{\mathbf{p}}^{r\dagger}e^{ip\cdot x} \Bigr),
  \label{eq:4.131}
\end{equation}
where $r = 0, 1, 2, 3$ labels a basis of polarization vectors.  The external
line factors in the Feynman rules above follow immediately from this
expansion, just as we obtained $u^s$'s and $v$'s as the external line factors
for Dirac particles.  The only subtlety is that we must restrict initial- and
final-state photons to be transversely polarized: Their polarization vectors
are always of the form $\epsilon^\mu = (0, \bm{\epsilon})$, where
$\mathbf{p}\cdot\bm{\epsilon} = 0$.  For $\mathbf{p}$ along the $z$-axis, the
right- and left-handed polarization vectors are
\begin{equation}
  \epsilon^{\mu}_{L/R} = (0, 1, \pm i, 0)/\sqrt{2} .
  \label{eq:4.132}
\end{equation}

The form of the QED vertex factor is also easy to justify, by simply looking
at the interaction Hamiltonian \eqref{eq:4.129}.  Note that the $\gamma^\mu$
matrix in a QED amplitude will sit between spinors or other $\gamma$ matrices,
with the Dirac indices contracted along the fermion line.  Note also that this
interaction term is specific to the case of an electron (and its antiparticle,
the positron).  In general, for a Dirac particle with electric charge
$Q|e|$,
\begin{equation}
  \text{vertex} = -iQ|e|\gamma^\mu .
  \label{eq:4.133}
\end{equation}
For example, an electron has $Q = -1$, an up quark has $Q = +2/3$, and a
down quark has $Q_d = -1/3$.

There is no easy way to derive the form of the photon propagator, so for now
we will settle for a plausibility argument.  Since the electromagnetic field
in Lorentz gauge obeys the massless Klein--Gordon equation, it should come as
no surprise that the photon propagator is nearly identical to the massless
Klein--Gordon propagator.  The factor of $-g^{\mu\nu}$, however, requires
explanation.  Lorentz invariance dictates that the photon propagator be an
isotropic second-rank tensor that can dot together the $\gamma^\mu$ and
$\gamma^\nu$ from the vertices at each end.  The simplest candidate is
$g^{\mu\nu}$.  To understand the overall sign of the propagator, evaluate its
Fourier transform:
\begin{equation}
  \int \frac{d^4q}{(2\pi)^4}\ \frac{-ig^{\mu\nu}}{q^2+i\epsilon}\
    e^{iq\cdot(x-y)}
  = \int \frac{d^3q}{(2\pi)^3}\ \frac{1}{2|\mathbf{q}|}\
    \Bigl( -g^{\mu\nu} \Bigr)\ \cdots ,
  \label{eq:4.132b}
\end{equation}
Presumably this is equal to $\langle 0|T[A^\mu(x)A^\nu(y)]|0\rangle$.  Now
set $\mu = \nu$, and take the limit $x^0\to y^0$ from the positive direction.
Then this quantity becomes the norm of the state $A^\mu(x)|0\rangle$, which
should be positive.  We see that our choice of signs in the propagator
implies that the three states created by $A^i$, with $i = 1, 2, 3$, indeed
have positive norm.  These states include all real (non-virtual) photons,
which always have spacelike polarizations.  Unfortunately, because $g^{\mu\nu}$
is not positive definite, the states created by $A^0$ inevitably have negative
norm.  This is potentially a serious problem for any theory with vector
particles.  For Quantum Electrodynamics, we will show in Section~5.5 that the
negative-norm states created by $A^0$ are never produced in physical
processes.  In Section~9.4 we will give a careful derivation of the photon
propagator.

The full set of Feynman rules for QED is:
\begin{enumerate}
  \item For each fermion line (with momentum $p$ flowing along the arrow),
    \begin{equation}
      \frac{i(\slashed{p}+m)}{p^2-m^2+i\epsilon} ;
    \end{equation}
  \item For each photon line,
    \begin{equation}
      \frac{-ig_{\mu\nu}}{q^2+i\epsilon} ;
    \end{equation}
  \item For each vertex, $-ie\gamma^\mu$;
  \item External fermions: initial fermions $u^s(p)$, initial antifermions
    $\bar{v}^s(p)$, final fermions $\bar{u}^s(p)$, final antifermions
    $v^s(p)$;
  \item External photons: $\epsilon^\mu(p)$ (initial), $\epsilon^{\mu*}(p)$
    (final);
  \item Impose momentum conservation at each vertex, integrate over
    undetermined loop momenta, and fix the overall sign.
\end{enumerate}

\subsection{The Coulomb Potential}

As a simple application of these Feynman rules, and to better understand the
sign of the propagator, let us repeat the nonrelativistic scattering
calculation of the previous section, this time for QED.  The leading-order
contribution is
\begin{equation}
  i\mathcal{M} = (-ie)^2\ \bar{u}(p')\gamma^\mu u(p)\
    \frac{-ig_{\mu\nu}}{q^2}\ \bar{u}(k')\gamma^\nu u(k) .
  \label{eq:4.133b}
\end{equation}
In the nonrelativistic limit,
\begin{equation}
  \bar{u}(p')\gamma^0 u(p) = u^{\dagger}(p')u(p) \approx +2m\,\xi'^{\dagger}\xi .
\end{equation}
You can easily verify that the other terms, $\bar{u}(p')\gamma^i u(p)$,
vanish if $\mathbf{p} = \mathbf{p}' = 0$; they can therefore be neglected
compared to $\bar{u}(p')\gamma^0 u(p)$ in the nonrelativistic limit.  Thus we
have
\begin{equation}
  \begin{aligned}
    i\mathcal{M} &\approx -e^2\ \frac{1}{|\mathbf{p}-\mathbf{p}'|^2}\
      (2m\,\xi'^{\dagger}\xi)_p\ (2m\,\xi'^{\dagger}\xi)_k\ \delta_{00}\\
    &= \frac{-e^2}{|\mathbf{p}-\mathbf{p}'|^2}\ 4m^2\ \cdots .
  \end{aligned}
  \label{eq:4.134}
\end{equation}
Comparing this to the Yukawa case, we see that there is an extra factor of
$-1$; the potential is a repulsive Yukawa potential with $m=0$, that is, a
repulsive Coulomb potential:
\begin{equation}
  V(\mathbf{r}) = \frac{e^2}{4\pi r} = \frac{\alpha}{r},
  \label{eq:4.135}
\end{equation}
where $\alpha = e^2/4\pi \approx 1/137$ is the fine-structure constant.

For particle-antiparticle scattering, note first that
\begin{equation}
  \bar{v}(k)\gamma^0 v(k') = v^{\dagger}(k)v(k') \approx +2m\,\eta'^{\dagger}\eta .
\end{equation}
The presence of the $\gamma^0$ eliminates the minus sign that we found in the
Yukawa case.  The nonrelativistic scattering amplitude is therefore
\begin{equation}
  i\mathcal{M} = \frac{+e^2}{|\mathbf{p}-\mathbf{p}'|^2}\ \cdots ,
\end{equation}
where the $(-1)$ is the same fermion minus sign we saw in the Yukawa case.
This is an attractive potential.  Similarly, for antifermion-antifermion
scattering one finds a repulsive potential.  We have just verified that in
quantum field theory, when a vector particle is exchanged, like charges repel
while unlike charges attract.

Note that the repulsion in fermion-fermion scattering came entirely from the
extra factor $-g_{00} = -1$ in the vector boson propagator.  A tensor boson,
such as the graviton, would have a propagator
\begin{equation}
  \frac{i}{q^2+i\epsilon}\ \Bigl[
    \tfrac{1}{2}(-g_{\mu\rho})(-g_{\nu\sigma}) + \tfrac{1}{2}(-g_{\mu\sigma})
      (-g_{\nu\rho}) - \tfrac{1}{3}(-g_{\mu\nu})(-g_{\rho\sigma}) \Bigr] ,
  \label{eq:4.136}
\end{equation}
which in nonrelativistic collisions gives a factor $(-g_{00})^2 = +1$; this
will result in a universally attractive potential.  It is reassuring to see
that quantum field theory does indeed reproduce the obvious features of the
electric and gravitational forces:
\begin{equation}
  \begin{array}{ccc}
    \text{Exchanged particle} & f\bar{f}\text{ and }\bar{f}\bar{f} & ff\\
    \hline
    \text{scalar (Yukawa)} & \text{attractive} & \text{attractive}\\
    \text{vector (electricity)} & \text{repulsive} & \text{attractive}\\
    \text{tensor (gravity)} & \text{attractive} & \text{attractive}
  \end{array}
  \label{eq:4.137}
\end{equation}

\problems

\textbf{4.1} Let us return to the problem of the creation of Klein--Gordon
particles by a classical source.  Recall from Chapter~2 that this process can
be described by the Hamiltonian
\begin{equation}
  H = H_0 + \int d^3x\ j(x)\phi(x) ,
\end{equation}
where $H_0$ is the free Klein--Gordon Hamiltonian, $\phi(x)$ is the
Klein--Gordon field, and $j(x)$ is a c-number scalar function.  We found
that, if the system is in the vacuum state before the source is turned on,
the source will create a mean number of particles
\begin{equation}
  \langle N\rangle = \int \frac{d^3p}{(2\pi)^3}\
    \frac{|\tilde{j}(p)|^2}{2E_{\mathbf{p}}} .
\end{equation}
In this problem we will verify that statement, and extract more detailed
information, by using a perturbation expansion in the strength of the source.
\begin{enumerate}[label=(\alph*)]
  \item Show that the probability that the source creates no particles is
    given by
    \begin{equation}
      P(0) = \Bigl| \langle 0|T\Bigl\{ \exp\Bigl[ i\int d^4x\
        j(x)\phi_I(x) \Bigr] \Bigr\}|0\rangle \Bigr|^2 .
    \end{equation}
  \item Evaluate the term in $P(0)$ of order $j^2$, and show that
    $P(0) = 1 - \lambda + \mathcal{O}(j^4)$, where $\lambda$ equals the
    expression given above for $\langle N\rangle$.
  \item Represent the term computed in part (b) as a Feynman diagram.  Now
    represent the whole perturbation series for $P(0)$ in terms of Feynman
    diagrams.  Show that this series exponentiates, so that it can be summed
    exactly: $P(0) = \exp(-\lambda)$.
  \item Compute the probability that the source creates one particle of
    momentum $\mathbf{k}$.  Perform this computation first to
    $\mathcal{O}(j)$ and then to all orders, using the trick of part (c) to
    sum the series.
  \item Show that the probability of producing $n$ particles is given by
    \begin{equation}
      P(n) = (1/n!)\,\lambda^n\exp(-\lambda) .
    \end{equation}
    This is a Poisson distribution.
  \item Prove the following facts about the Poisson distribution:
    \begin{equation}
      \sum_{n=0}^{\infty}P(n) = 1;\qquad
      \langle N\rangle = \sum_{n=0}^{\infty} n\,P(n) = \lambda .
    \end{equation}
    The first identity says that the $P(n)$'s are properly normalized
    probabilities, while the second confirms our proposal for $\langle N\rangle$.
    Compute the mean square fluctuation
    $\bigl( \Delta N \bigr)^2 = \langle (N - \langle N\rangle)^2 \rangle$.
\end{enumerate}

\textbf{4.2} \emph{Decay of a scalar particle.}  Consider the following
Lagrangian, involving two real scalar fields $\Phi$ and $\phi$:
\begin{equation}
  \mathcal{L} = \tfrac{1}{2}(\partial_\mu\Phi)^2 - \tfrac{1}{2}M^2\Phi^2
    + \tfrac{1}{2}(\partial_\mu\phi)^2 - \tfrac{1}{2}m^2\phi^2
    - \mu\Phi\phi\phi .
\end{equation}
The last term is an interaction that allows a $\Phi$ particle to decay into
two $\phi$'s, provided that $M > 2m$.  Assuming that this condition is met,
calculate the lifetime of the $\Phi$ to lowest order in $\mu$.

\textbf{4.3} \emph{Linear sigma model.}  The interactions of pions at low
energy can be described by a phenomenological model called the linear sigma
model.  Essentially, this model consists of $N$ real scalar fields coupled by
a $\phi^4$ interaction that is symmetric under rotations of the $N$ fields.
More specifically, let $\Phi^i(x)$, $i = 1, \ldots, N$ be a set of $N$ fields,
governed by the Hamiltonian
\begin{equation}
  H = \int d^3x\ \Bigl[ \tfrac{1}{2}\pi^i\pi^i
    + \tfrac{1}{2}\nabla\Phi^i\cdot\nabla\Phi^i
    + V(\Phi^2) \Bigr],
\end{equation}
where $(\Phi^2) = \Phi\cdot\Phi$, and
\begin{equation}
  V(\Phi^2) = \tfrac{1}{2}m^2(\Phi^2) + \tfrac{\lambda}{4!}(\Phi^2)^2
\end{equation}
is a function symmetric under rotations of $\Phi$.  For (classical) field
configurations of $\Phi^2(x)$ that are constant in space and time, this term
gives the only contribution to $H$; hence, $V$ is the field potential energy.

(What does this Hamiltonian have to do with the strong interactions?  There
are two types of light quarks, $u$ and $d$.  These quarks have identical
strong interactions, but different masses.  If these quarks are massless, the
Hamiltonian of the strong interactions is invariant to unitary
transformations of the 2-component object $(u,d)$:
\begin{equation}
  \begin{pmatrix}u\\ d\end{pmatrix} \to U\begin{pmatrix}u\\ d\end{pmatrix} .
\end{equation}
This transformation is called an isospin rotation.  If, in addition, the
strong interactions are described by a vector ``gluon'' field (as is true in
QCD), the strong interaction Hamiltonian is invariant to the isospin rotations
done separately on the left-handed and right-handed components of the quark
fields.  Thus, the complete symmetry of QCD with two massless quarks is
$SU(2)\times SU(2)$.  It happens that $SO(4)$, the group of rotations in 4
dimensions, is isomorphic to $SU(2)\times SU(2)$, so for $N=4$, the linear
sigma model has the same symmetry group as the strong interactions.)
\begin{enumerate}[label=(\alph*)]
  \item Analyze the linear sigma model for $m^2>0$ by noticing that, for
    $\lambda=0$, the Hamiltonian given above is exactly $N$ copies of the
    Klein--Gordon Hamiltonian.  We can then calculate scattering amplitudes
    as perturbation series in the parameter $\lambda$.  Show that the
    propagator is
    \begin{equation}
      \Phi^i(x)\Phi^j(y) = \delta^{ij}D_F(x-y) ,
    \end{equation}
    where $D_F$ is the standard Klein--Gordon propagator for mass $m$, and
    that there is one type of vertex given by
    \begin{equation}
      -2i\lambda\bigl( \delta^{ij}\delta^{kl} + \delta^{il}\delta^{jk}
        + \delta^{ik}\delta^{jl} \bigr) .
      \label{eq:4.vertex}
    \end{equation}
    (That is, the vertex between two $\Phi^1$'s and two $\Phi^2$'s has the
    value $(-2i\lambda)$; that between four $\Phi^1$'s has the value
    $(-6i\lambda)$.)  Compute, to leading order in $\lambda$, the
    differential cross sections $d\sigma/d\Omega$, in the center-of-mass
    frame, for the scattering processes
    \begin{equation}
      \Phi^1\Phi^2 \to \Phi^1\Phi^2,
      \qquad \Phi^1\Phi^1 \to \Phi^2\Phi^2,
      \qquad \Phi^1\Phi^1 \to \Phi^1\Phi^1 .
    \end{equation}
    Write your answers in terms of the Mandelstam variables $s$, $t$, $u$.
  \item For $m^2<0$, the potential $V$ has a minimum at a nonvanishing value
    of $|\Phi|$, given by
    \begin{equation}
      |\Phi_0|^2 = v^2 = -\frac{6m^2}{\lambda} .
    \end{equation}
    At this value of the field, the $SO(N)$ symmetry is spontaneously broken
    down to $SO(N-1)$.  To study this theory, let $\Phi$ vary only in the
    vicinity of its minimum, and define a new field $\pi^i(x)$ by
    \begin{equation}
      \Phi^i(x) = \begin{cases}
        v + \sigma(x) & i = N\\
        \pi^i(x) & i \neq N
      \end{cases}
    \end{equation}
    Rewrite the Lagrangian in terms of $\pi$ and $\sigma$.  (There will be
    $\sigma\pi\pi$ and $\sigma\sigma\sigma$ and $\sigma\sigma\sigma\sigma$
    and $\sigma\sigma\pi\pi$ and $\sigma\pi\pi\pi$ vertices.)
  \item For $m^2 < 0$, analyze the theory near the minimum of $V$.  Show
    that the $N-1$ fields $\pi^i$ are massless (these are the Nambu--Goldstone
    bosons, which we will discuss in Chapter~11), while the $\sigma$ field
    has mass
    \begin{equation}
      m_\sigma^2 = -\frac{2}{3}m^2 .
    \end{equation}
    Compute, to leading order in $\lambda$, the differential cross section
    $d\sigma/d\Omega$, in the center-of-mass frame, for the scattering
    process
    \begin{equation}
      \pi^1\pi^2 \to \pi^1\pi^2 .
    \end{equation}
    There are now four Feynman diagrams that contribute.  Show that, at
    threshold ($p^\mu = 0$), these diagrams sum to zero.  (\emph{Hint}: It
    may be easiest to first consider the specific process
    $\pi^1\pi^1\to\pi^2\pi^2$, for which only the first and fourth diagrams
    are nonzero, before tackling the general case.)  Show that, in the
    special case $N=2$ (1 species of pion), the term of $\mathcal{O}(p^2)$
    also cancels.
  \item Add to $V$ a symmetry-breaking term,
    \begin{equation}
      \Delta V = -a\Phi^N ,
    \end{equation}
    where $a$ is a (small) constant.  (In QCD, a term of this form is
    produced if the $u$ and $d$ quarks have the same nonvanishing mass.)
    Find the new value of $v$ that minimizes $V$, and work out the content
    of the theory about that point.  Show that the pion acquires a mass such
    that $m_\pi \sim a$, and show that the pion scattering amplitude at
    threshold is now nonvanishing and also proportional to $a$.
\end{enumerate}

\textbf{4.4} \emph{Rutherford scattering.}  The cross section for scattering
of an electron by the Coulomb field of a nucleus can be computed, to lowest
order, without quantizing the electromagnetic field.  Instead, treat the
field as a given, classical potential $A^\mu(x)$.  The interaction Hamiltonian
is
\begin{equation}
  H_{\mathrm{int}} = -e\int d^3x\ \bar{\psi}\gamma^\mu\psi\,A_\mu ,
\end{equation}
where $\psi(x)$ is the usual quantized Dirac field.
\begin{enumerate}[label=(\alph*)]
  \item Show that the $T$-matrix element for electron scattering off a
    localized classical potential is, to lowest order,
    \begin{equation}
      \langle \mathbf{p}'|iT|\mathbf{p}\rangle
        = -ie\ \bar{u}(p')\gamma^\mu u(p)\ \tilde{A}_\mu(p'-p) ,
    \end{equation}
    where $\tilde{A}_\mu(q)$ is the four-dimensional Fourier transform of
    $A_\mu(x)$.
  \item If $A_\mu(x)$ is time independent, its Fourier transform contains a
    delta function of energy.  It is then natural to define
    \begin{equation}
      \langle \mathbf{p}'|iT|\mathbf{p}\rangle
        = i\mathcal{M}\cdot(2\pi)\,\delta(E_f - E_i) ,
    \end{equation}
    where $E_i$ and $E_f$ are the initial and final energies of the
    particle, and to adopt a new Feynman rule for computing $\mathcal{M}$:
    \begin{equation}
      \text{vertex} = -ie\gamma^\mu\ \tilde{A}_\mu(\mathbf{q}),
    \end{equation}
    where $\tilde{A}_\mu(\mathbf{q})$ is the three-dimensional Fourier
    transform of $A_\mu(\mathbf{x})$.  Given this definition of $\mathcal{M}$,
    show that the cross section for scattering off a time-independent,
    localized potential is
    \begin{equation}
      \frac{d\sigma}{d\Omega} = \frac{1}{4\pi^2}\ |\mathcal{M}
        (\mathbf{p}'\to\mathbf{p})|^2\ \frac{p'^2\,dp'/dE'}{(2\pi)^3\,E'}
        \cdots ,
    \end{equation}
    a natural modification of \eqref{eq:4.79}.  Integrate over $|p_f|$ to
    find a simple expression for $d\sigma/d\Omega$.
  \item Specialize to the case of electron scattering from a Coulomb
    potential ($A^0 = Ze/4\pi r$).  Working in the nonrelativistic limit,
    derive the Rutherford formula,
    \begin{equation}
      \frac{d\sigma}{d\Omega} = \frac{\alpha^2Z^2}
        {4m^2v^4\sin^4(\theta/2)} .
    \end{equation}
    (With a few calculational tricks from Section~5.1, you will have no
    difficulty evaluating the general cross section in the relativistic case;
    see Problem~5.1.)
\end{enumerate}
